\chapter{纤维丛}

本章介绍矢量丛、标架丛及其联络，并概述纤维丛与规范场的关系。
沿用前文记法：矢量、对偶矢量和张量分别记作 $v$、$\omega$、$T$；对偶配对为 $\omega(v)$。
选定基或局部标架后，写作 $v=v^i e_i$、$\omega=\omega_i e^i$，其中 $e^i(e_j)=\delta^i_j$。
分量指标均相对于所选基，丛标架指标与底流形坐标指标分别取值。

矢量丛的定义不依赖 Lie 群。\S\ref{chfb:sec_fb}、\S\ref{chfb:sec_frame-bundles}只需 Lie 群作用的概念；
\S\ref{chfb:sec_G-connection}还需 Lie 代数知识。

\section{丛}
丛是矢量丛和纤维丛的基本结构。

\begin{definition}
    \textbf{丛}（bundle）是一个三元组$(E, B, \pi)$，其中$\pi: E\to B$（或$E\xrightarrow{\pi} B$）是一个连续满射；
    拓扑空间$B$称为\textbf{底空间}，拓扑空间$E$称为\textbf{总空间}或 \textbf{丛空间} ，映射$\pi$称为\textbf{丛投影}；
    $\forall b\in B$，空间$\pi^{-1}(b)\subset E$称为植于$ b\in B$点的\textbf{纤维}（fibre）。
\end{definition}
直观地，我们认为丛是纤维$\pi^{-1}(b)$的并集，并被空间$E$的拓扑“粘在一起”。

\begin{definition}
     如果$E'$ 是 $E$ 的子空间，$B'$ 是 $B$ 的子空间，且 $ \pi' = \pi |_{E'}: E'\to B'$；
     那么称丛 $(E',  B',\pi')$ 是 $(E, B,\pi)$ 的\textbf{子丛}。
\end{definition}

\begin{definition}
    若映射$s: B\to E$使得$\pi \circ s = \id $成立，则称$s$是丛$(E, B,\pi)$的\textbf{截面}。
    换句话说，截面$s$把$b\in B$点映射成纤维$\pi^{-1}(b)$中的元素，即$s(b)\in \pi^{-1}(b)$。
\end{definition}

若$(E',B,\pi|_{E'})$是同一底空间上的子丛，则截面$s$也是该子丛的截面，当且仅当$s(B)\subset E'$。

设 $B$、$F$ 为拓扑空间，$\pi:B\times F\to B$ 为第一因子投影，则 $(B\times F,B,\pi)$ 构成积丛。
其截面恰为 $s(b)=(b,f(b))$，其中 $f:B\to F$ 由 $s$ 唯一确定。
事实上，若 $s(b)=(\tau(b),f(b))$，则截面条件 $\pi\circ s=\id$ 等价于 $\tau(b)=b$。

因此，积丛的截面与映射$f:B\to F$一一对应。
若$E\subset B\times F$是同底空间的子丛，只须再要求$(b,f(b))\in E$。

截面在底空间的每点选取一个纤维元素；若只沿一条曲线选取，就得到曲线的提升。
\begin{definition}\label{chfb:def_curve-lift}
    设$\pi:E\to B$为丛，$\gamma:I\to B$为曲线，$I\subset\R$为区间。
    若曲线$\widetilde\gamma:I\to E$满足
    \eq{\pi\circ\widetilde\gamma=\gamma,}
    则称$\widetilde\gamma$为$\gamma$的\textbf{提升}(lift of a curve)。
    即$\widetilde\gamma(t)\in\pi^{-1}(\gamma(t))$，投影后还原原曲线。此处曲线均连续；讨论光滑丛时，所论曲线及提升均取光滑。
\end{definition}

例如，积丛中的提升具有形式$\widetilde\gamma(t)=(\gamma(t),f(t))$，其中$f:I\to F$可任取。
故提升本身不需要联络，指定初点后通常也不唯一。
若$s$为截面，则$s\circ\gamma$是一个提升；但一般提升未必来自底空间上的截面，因为当$\gamma$回到同一点时，提升可以到达该纤维中的另一点。

\section{矢量丛定义}

从\S \ref{chdm:sec_tangent-bundles}的微分流形$TM$所具有的特殊构造可抽象出如下定义：
\begin{definition}\label{chfb:def_vector-bundles}
    设$E$、$M$是两个光滑流形，维数分别是$q+m$和$m$；$\pi:E\to M$是光滑满映射；
    $V := \mathbb{R}^q$是$q$维矢量空间。假设流形$M$上存在一个开覆
    盖$\{U_\alpha; \alpha\in I\}$以及一组映射$\{\psi_\alpha; \alpha\in I\}$；
    $\pi^{-1}(U_\alpha)$是流形$E$的开覆盖；
    它们满足如下条件：

    (1) $\forall \alpha \in I$，映射$\psi_\alpha$是从$\pi^{-1}(U_\alpha)$到
    $U_\alpha \times \mathbb{R}^q$的微分同胚，且$\forall p \in U_\alpha$有
    \begin{align*}
        \psi_\alpha\left(v\right) := \left(p, y \right),
        \quad \ v\in \pi^{-1} (U_\alpha),\ y \in \mathbb{R}^q ;
        \quad \text{且}\quad
        \pi \circ \psi_\alpha^{-1} (p,y) = p.
    \end{align*}

    (2) 对于任意固定的$p\in U_\alpha$，$\psi^{-1}_{\alpha,p}(\cdot) := \psi^{-1}_\alpha(p,\cdot):\mathbb{R}^q\to \pi^{-1}(p)$是
    微分同胚映射；而且当$p\in U_\alpha \cap U_\beta \neq \varnothing$时，
    映射（称为\textbf{过渡函数族}或转移函数族）
    \begin{equation*}
        g^p_{\beta\alpha} (\cdot) := \psi_{\beta,p} \circ \psi^{-1}_{\alpha,p}(\cdot ) : \mathbb{R}^q\to\mathbb{R}^q ,
        \qquad g^p_{\beta\alpha}\text{的上角标$p$可省略}
    \end{equation*}
    是线性同构，也就是说$g^p_{\beta\alpha}\in GL(q,\mathbb{R})$，$GL(q,\mathbb{R})$是{实}一般线性群。

    (3) 当$U_\alpha \cap U_\beta \neq \varnothing$时，
    映射$g_{\beta\alpha}: U_\alpha \cap U_\beta \to GL(q,\mathbb{R})$是$C^\infty$的。

    则称$(E,M,\pi)$是光滑流形$M$上秩为$q$的实\textbf{矢量丛}，其中$E$称为\textbf{丛空间}，
    $M$称为\textbf{底空间}（底流形），$\pi$称为\textbf{丛投影}，$V=\mathbb{R}^q$称为\textbf{纤维型}。
\qed
\end{definition}

\begin{definition}
把定义\ref{chfb:def_vector-bundles}中的实空间$\mathbb{R}^q$换成{复}矢量空间$\mathbb{C}^q$，
并把$GL(q,\mathbb{R})$换成复群$GL(q,\mathbb{C})$；则称$\bigl(E,M,\pi,\mathbb{C}^q,GL(q,\mathbb{C})\bigr)$是
流形$M$上的\textbf{复矢量丛}。\qed
\end{definition}

条件(1)中的映射$\psi_\alpha:\pi^{-1}(U_\alpha) \to U_\alpha \times \mathbb{R}^q$称为
矢量丛$E$在$U_\alpha$上的\textbf{局部平凡化}，
$U_\alpha \times \mathbb{R}^q (\ \xleftrightarrow[\varphi^{-1}]{\varphi}
\mathbb{R}^m \times \mathbb{R}^q\ )$也就是
纤维$\pi^{-1}(U_\alpha)$的局部坐标卡。

由定义\ref{chfb:def_vector-bundles}中条件(2)可知，
若映射$\psi^{-1}_{\alpha,p}(\cdot)$、$\psi^{-1}_{\beta,p}(\cdot)$分别
将$\mathbb{R}^q$中的两个元素$y_\alpha$、$y_\beta$映射成丛$E$中的同一个点，则有
\begin{equation}\label{chfb:eqn_gbayy}
    \psi^{-1}_{\alpha,p}(y_\alpha) = \psi^{-1}_{\beta,p}(y_\beta)
    \quad \Leftrightarrow \quad
    g^p_{\beta \alpha}(y_\alpha) = y_\beta .
\end{equation}
我们把$\mathbb{R}^q$中的元素记成列矢量，即$y=(y^1,\cdots, y^q)^{\mathrm{T}}$；
那么上式可以理解成：$GL(q)$是左作用在$\mathbb{R}^q$上的Lie 变换群。
有了这样的理解，很容易将矢量丛概念推广到一般纤维丛。

矢量丛定义\ref{chfb:def_vector-bundles}中的条件(1)和(2)结合起来说明，
丛空间$E$可以看成一堆笛卡尔积空间$\{U_\alpha \times \mathbb{R}^q\}$；
并且当$U_\alpha \cap U_\beta \neq \varnothing$时，$\forall p\in U_\alpha \cap U_\beta$，
应该将纤维$\{p\}\times \mathbb{R}^q \subset U_\alpha \times \mathbb{R}^q$和
纤维$\{p\}\times \mathbb{R}^q \subset U_\beta \times \mathbb{R}^q$粘合起来，
粘合规则就是\eqref{chfb:eqn_gbayy}式。

由下面命题可见：矢量丛 $\pi: E \to M$在直观上可以想象为“植”在光滑流形$M$上的一“簇”矢量空间；
任一点$p\in M$的纤维$\pi^{-1}(p)$都是一个线性空间。

\begin{proposition}\label{chfb:thm_linear-space}
设 $\pi: E \rightarrow M$是光滑流形$M$上秩为$q$的矢量丛，
则在每一点$p \in M$上的纤维 $\pi^{-1}(p)$有自然的线性结构，使它成为 $q$ 维矢量空间。
\end{proposition}

\begin{proof}
设 $p \in U_\alpha$，$\psi_\alpha^{-1}: U_\alpha \times \mathbb{R}^q \rightarrow \pi^{-1} (U_\alpha)$是
局部平凡化映射。在纤维$\pi^{-1}(p)$中定义加法和数乘法如下：
对于任意的 $v, w \in \pi^{-1}(p)$ 及 $\lambda \in \mathbb{R}$，令
\begin{equation}\label{chfb:eqn_tmp-vlwt}
    v+\lambda \cdot w := \psi_{\alpha, p}^{-1} \bigl(\psi_{\alpha, p} (v)
    +\lambda \cdot \psi_{\alpha, p} (w)\bigr) .
\end{equation}
我们已知$\psi_{\alpha, p} (v)\in \mathbb{R}^q$，故$\psi_{\alpha, p} (v)$自身是线性空间$\mathbb{R}^q$中
的矢量；我们借助$\psi_{\alpha, p}$、$\psi_{\alpha, p}^{-1}$是同胚映射，将这种线性结构传递至$\pi^{-1}(p)$中。
由式\eqref{chfb:eqn_tmp-vlwt}定义的$v+ \lambda w $是纤维$ \pi^{-1}(p)$中的元素。

我们先证明：在 $\pi^{-1}(p)$中由式\eqref{chfb:eqn_tmp-vlwt}定义的加法和数乘与
局部平凡化$\psi_\alpha$的选取无关。
实际上，若有 $p \in U_\alpha \cap U_\beta,\ \psi_\beta: \pi^{-1} (U_\beta )
 \rightarrow U_\beta \times \mathbb{R}^q $是另一个局部平凡化映射，
则由定义\ref{chfb:def_vector-bundles}中的条件(2)可知，
存在过渡函数$g_{\beta \alpha}^p \in GL(q)$使得
$g_{\beta \alpha}^p =\psi_{\beta, p} \circ \psi_{\alpha, p}^{-1}: \mathbb{R}^q \rightarrow \mathbb{R}^q $，
并且$g_{\beta \alpha}^p$是左作用在$\mathbb{R}^q$上的线性变换。因此
$g_{\beta \alpha}^p \circ \psi_{\alpha, p}=\psi_{\beta, p}: \pi^{-1}(p) \rightarrow \mathbb{R}^q $。
因$g_{\beta \alpha}^p$是线性变换，故有
$    \psi_{\beta, p}(v )+\lambda \cdot \psi_{\beta, p} (w)=
    g_{\beta \alpha}^p \cdot \bigl(\psi_{\alpha, p}(v)+ \lambda \psi_{\alpha, p} (w)\bigr) $。
由此式可得

\begin{align*}
    \psi_{\beta, p}^{-1} \bigl(\psi_{\beta, p} (v)+\lambda \psi_{\beta, p} (w)\bigr)
    = \psi_{\beta, p}^{-1} \circ g_{\beta \alpha}^p \bigl(\psi_{\alpha, p}(v)+\lambda \psi_{\alpha, p}(w)\bigr)
    = \psi_{\alpha, p}^{-1} \bigl(\psi_{\alpha, p} (v)+\lambda \psi_{\alpha, p} (w)\bigr) .
\end{align*}
上式说明了：式\eqref{chfb:eqn_tmp-vlwt}定义的加法和数乘与局部平凡化的选取无关。

接着，容易验证$\pi^{-1}(p)$关于式\eqref{chfb:eqn_tmp-vlwt}定义的加法和数乘
满足线性空间的八条公理（留给读者），零元素是 $\psi_{\alpha, p}^{-1}(0)$。
\end{proof}

在把$E$看成一堆积流形 $\left\{U_\alpha \times \mathbb{R}^q\right\}$并且同一点上的
各条纤维保持线性关系地粘合起来的结果时，过渡函数族$\left\{g_{\alpha \beta}: U_\alpha \cap U_\beta
\rightarrow GL(q)\right\}$起到核心的作用；它满足如下条件。

\begin{proposition}\label{chfb:thm_trans-func}
(1) $g_{\alpha \alpha}^p=I,\ \forall p \in U_\alpha$，$I$是 $q \times q$ 单位矩阵；

(2) 若$U_\alpha \cap U_\beta \cap U_\gamma \neq \varnothing$，
则$\forall p \in U_\alpha \cap U_\beta \cap U_\gamma$ 有
$g_{\alpha \beta}^p \cdot g_{\beta \gamma}^p=g_{\alpha \gamma}^p $；

(3) 对于 $p \in U_\alpha \cap U_\beta$ 有$g_{\alpha \beta}^p=(g_{\beta \alpha}^p)^{-1} $。
\end{proposition}

\begin{proof}
    (1)： 由定义，有
    $ g_{\alpha \alpha}^p=\psi_{\alpha, p} \circ \psi_{\alpha, p}^{-1} = \id     $，
    自然对应$q \times q$ 单位矩阵。

(2)： 由定义，有$g_{\alpha \beta}^p=\psi_{\alpha, p} \circ \psi_{\beta, p}^{-1}$，
$g_{\beta \gamma}^p=\psi_{\beta, p} \circ \psi_{\gamma, p}^{-1}$。
所以当 $p \in U_\alpha \cap U_\beta \cap U_\gamma$ 时，有
$  g_{\alpha \beta}^p \cdot g_{\beta \gamma}^p
    =  g_{\alpha \beta}^p \circ g_{\beta \gamma}^p
    =  \psi_{\alpha, p} \circ \psi_{\gamma, p}^{-1}=g_{\alpha \gamma}^p$。

(3)： 在条件 (2) 中令 $\gamma=\alpha$，则
$g_{\alpha \beta}^p \cdot g_{\beta \alpha}^p=g_{\alpha\alpha}^p=I$，
因此$g_{\alpha \beta}^p=\left(g_{\beta \alpha}^p\right)^{-1} $。

这个命题对一般纤维丛的过渡函数族同样成立。
\end{proof}

\begin{theorem}\label{chfb:thm_vector-existence}
设 $M$ 是一个 $m$ 维光滑流形，$\left\{U_\alpha: \alpha \in I\right\}$是 $M$ 的一个开覆盖。
如果对于任意一对指标 $\alpha, \beta \in I$，在 $U_\alpha \cap U_\beta \neq \varnothing$ 时，
都指定了一个光滑映射 $g_{\alpha \beta}: U_\alpha \cap U_\beta \rightarrow GL(q)$，
使得它们满足命题\ref{chfb:thm_trans-func}中的条件 (1)和 (2)；
那么存在一个秩为 $q$ 的矢量丛$(E,M,\pi)$，
以 $\left\{g_{\alpha \beta}: U_\alpha \cap U_\beta \rightarrow GL(q)\right\}$ 为它的过渡函数族。
\end{theorem}

\begin{proof}
    见定理\ref{chfb:thm_fb-existence}。
    本定理说明过渡函数族是构造矢量丛的核心。
\end{proof}

\begin{definition}
    设$(E,M,\pi)$是矢量丛，若$s:M\to E$是一个$C^\infty$映射，
    并且满足$\pi \circ s = \id  :M\to M$，则称映射$s$是矢量丛$E$的一个\textbf{截面}。
    若$s:U\to E$，其中$U$是$M$某坐标卡中的开集，则称映射$s$是定义在$U$上的\textbf{局部截面}。
    光滑截面的集合记为$\Gamma(E)$。
\end{definition}

\begin{proposition}
    设$\pi: E \rightarrow M$是一个矢量丛，则截面空间 $\Gamma(E)$ 是一个 $C^{\infty}(M)$ - 模。
\end{proposition}

\begin{proof}
    命题\ref{chfb:thm_linear-space}中已经证明矢量丛的每条纤维$\pi^{-1}(p)$是
    与纤维型$V = \mathbb{R}^q$同构的线性空间；故可以逐点定义截面的加法和
    实值函数的乘法。设$s,t\in \Gamma(E)$，$f$是$M$上的实值光滑函数，则$\forall p\in M$，定义
    \begin{equation}
        (s+t)(p):= s(p) + t(p);\qquad
        (fs)(p) := f(p) \cdot s(p) .
    \end{equation}
    则$s+t$、$fs$仍是矢量丛$E$的光滑截面。这便证明了$\Gamma(E)$ 是 $C^{\infty}(M)$ - 模；
    $\Gamma(E)$自然也是实线性空间，通常是无穷维的。
\end{proof}

设矢量丛$(E,M,\pi)$的秩为$q$，$\{U_\alpha\}$是$M$的开覆盖。用$\{e_i\}$表示$\mathbb{R}^q$的一组基底，
则由局部平凡化映射，有
\begin{equation}
    s_i^{(\alpha)}(p) = \psi^{-1}_{\alpha}(p, e_i),\qquad \forall p\in U_\alpha,
    \quad \alpha \text{固定}, \quad 1 \leqslant i \leqslant q.
\end{equation}
便可得到$q$个映射$s_i^{(\alpha)} : U_\alpha \to E$（$\alpha$固定）。
显然有$\pi \circ s_i^{(\alpha)} =\id $。
这说明$\{s_i^{(\alpha)}(p)\}$是$q$个矢量丛$E$上的局部截面；
由于$\{e_i\}$是线性无关的，则$\{s_i^{(\alpha)}(p)\}$也是线性无关的。
$\{s_i^{(\alpha)}(p)\}$称为定义在$U_\alpha$上矢量丛$E$的\textbf{局部标架场}。

上述描述方式说明了$M$上的矢量丛$E$上一定存在$q$个局部标架场，
但整个$M$上未必存在这样的标架场。

\section{矢量丛例子}

\begin{example}\label{chfb:exm_tb}
    按照定义\ref{chfb:def_vector-bundles}，\S \ref{chdm:sec_tangent-bundles}中叙述的$\pi:TM\to M$是$M$上
    秩为$m$的矢量丛，称为微分流形$M$上的\textbf{切矢量丛}，简称\textbf{切丛}。
\end{example}
在\S \ref{chdm:sec_tangent-bundles}已经给出了切丛的微分结构；这个微分结构的局部坐标
表达式为$(p,v)$，其中$p$是底流形$M$上任意一点，$v$是$T_pM$中任一矢量。
我们定义丛投影为$\pi\bigl( (p,v)\bigr)=p$。
由\S \ref{chdm:sec_tangent-bundles}式\eqref{chdm:eqn_xy-transform}
\begin{equation}
    v^i = \tilde{v}^j \left. \frac{\partial {x}^i_\alpha}{\partial {x}^j_\beta}  \right | _p
    {\quad \Leftrightarrow \quad}
    \tilde{v}^j = {v}^i \left. \frac{\partial {x}^j_\beta}{\partial {x}^i_\alpha}  \right | _p
    \tag{\ref{chdm:eqn_xy-transform}}
\end{equation}
可知，过渡函数族$g_{\beta\alpha}: U_\alpha \cap U_\beta \to GL(q)$是
局部坐标卡变换的Jacobi矩阵：
\begin{equation}
    g_{\beta\alpha} = \psi_\beta \circ \psi^{-1}_{\alpha}
    = \frac{\partial {x}^j_\beta}{\partial {x}^i_\alpha}  .
    \qquad \text{注意}g_{\beta\alpha}\text{是} m\times m \text{阶矩阵}
\end{equation}

切丛截面在每点 $p$ 指定一个切矢 $v_p\in T_pM$，故切丛的光滑截面就是光滑切矢量场。
\qed

\begin{example}
    设流形$M$是单位圆周，试表示其切丛。
\end{example}
$M$的局部坐标为$\{x\}$取为
\begin{equation*}
    x=\left\{\begin{array}{ccc}
        \theta& -\epsilon <\theta<\pi+\epsilon & \text{开集记为}U \\
        \phi & \pi -\epsilon <\phi <2\pi +\epsilon  & \text{开集记为}V
    \end{array} \right.
    \qquad \text{其中} 0<\epsilon<\frac{\pi}{2}
\end{equation*}
当$U\cap V\neq \varnothing$时，两个坐标有$C^\infty$相容性关系：$\phi=\theta + \pi$。
容易算出此时的过渡函数族是$g_{UV}=g_{VU}=1$。
圆周上任一点的切矢量是$c'(x)=(-\sin x,\ \cos x)$。
$\forall p=(\cos x,\ \sin x)\in M$，令$v_p=(-\sin x,\ \cos x)_p$。
那么我们可以定义单位圆周切丛的丛映射为$\pi \bigl( (p,v_p) \bigr)=p$，
此切丛的局部平凡化是$\psi\bigl(\pi^{-1}(p) \bigr)=(p,v_p)$。
\qed

\begin{example}\label{chfb:exm_ctb}
    余切丛$T^* M$。
\end{example}
与切丛对偶，自然存在余切丛；我们较为详尽地描述它。
设$M$是$m$维光滑流形，$p\in M$点的余切空间记为$T_p^* M$，记
\begin{equation}
    T^* M := \bigcup _{p\in M} T_p^* M = \{\omega \in T_p^* M \mid p\in M\} .
\end{equation}
$T^*M$是流形$M$上全体余切矢量的集合。我们要给它赋予一个拓扑结构和微分结构，使它成为一个$2m$维的光滑流形。
与切丛情形类似，我们将略去如下证明：
$T^*M$是Hausdorff空间；$T^*M$存在开覆盖$\bigcup _{\alpha\in I} \pi^{-1} (U_\alpha)$，
并且这族开覆盖构成$T^*M$的可数拓扑基。

下面建立$T^*M$拓扑。定义丛投影映射$\pi:T^*M\to M$如下：
\begin{equation*}
    \pi(\omega) = p; \qquad \forall \omega \in T_p^* M,\quad p\in M.
\end{equation*}
这样，对每一个点$p\in M$，都有$\pi^{-1}(p)=T_p^* M$；容易看出映射$\pi$是满射。

设微分流形$M$有微分结构$\mathscr{A}=\{(U_\alpha,\varphi_\alpha;x^i)|\alpha \in I \}$，令
\begin{equation*}
    \pi^{-1} (U_\alpha) = \bigcup _{p\in U_\alpha} T_p^*M .
\end{equation*}
当$\alpha$遍历整个指标集$I$后，
便有$\bigcup _{\alpha \in I}  \pi^{-1} (U_\alpha) = T^*M$。
$\pi^{-1} (U_\alpha)$是$T^*M$的子集；
因$T^*_pM$是$m$维的，而点$p\in U_\alpha$（$U_\alpha$是$m$维）又是可变
动的，所以可知$\pi^{-1} (U_\alpha)$是$2m$维；进而$T^*M$也是$2m$维的。
这个子集中的元素自然是余切矢量。

设$(x_p^1,\cdots,x_p^m)$是$p$点在映射$\varphi_\alpha$作用下的坐标；
矢量$\omega \in \pi^{-1} (U_\alpha)$可以在$p$点自然余切
基矢$\left.\d x^i \right | _p$下展开，
其分量是$(\omega_1,\cdots,\omega_m)$。对于任意$\alpha \in I$，
定义局部平凡化映射$\psi_\alpha: \pi^{-1} (U_\alpha) \to U_\alpha \times \mathbb{R}^m
\overset{\varphi_\alpha}{\longrightarrow} \varphi_\alpha(U_\alpha) \times \mathbb{R}^m $
（其中$\varphi_\alpha(U_\alpha) \subset \mathbb{R}^m$维数是$m$）。
\begin{equation}\label{chfb:eqn_ctb_atalas}
    \psi_\alpha\left(\omega_i \left. \d x^i \right| _p \right) :=
    \left(x_p^1,\cdots,x_p^m, \omega_1,\cdots,\omega_m\right),
    \ \forall p\in U_\alpha,\ \omega\in \pi^{-1} (p).
\end{equation}
因$(x_p^1,\cdots,x_p^m)$与$(\omega_1,\cdots,\omega_m)$都是可变动的，所以上述
定义更好地说明了$\pi^{-1} (U_\alpha)$是$2m$维的。
可见映射$\psi_\alpha$是双射，其逆存在；
我们可以规定$\psi_\alpha$和$\psi_\alpha^{-1}$是连续映射，
相当于规定了$\pi^{-1} (U_\alpha)$中的开集；这样，
式\eqref{chfb:eqn_ctb_atalas}中的映射便是拓扑同胚的，
它将$\varphi_\alpha(U_\alpha) \times \mathbb{R}^m $中的（通常）拓扑
携带到$\pi^{-1} (U_\alpha)$中。
$(\pi^{-1} (U_\alpha),\psi_\alpha)$是$TM$一个坐标卡，
这样我们建立了$T^*M$的拓扑结构。

我们开始建立$T^*M$上的微分结构；只需验证坐标卡$\bigl(\pi^{-1} (U_\alpha),\psi_\alpha\bigr)$是$C^\infty$相容的。
实际上，$\pi^{-1} (U_\alpha) \cap \pi^{-1} (U_\beta) \neq \varnothing \ \Leftrightarrow\
U_\alpha \cap U_\beta \neq \varnothing$。
当$U_\alpha \cap U_\beta \neq \varnothing$时，设分别存在
坐标卡$(U_\alpha;x^i_\alpha)$和$(U_\beta;x^i_\beta)$；则，它们间的坐标变换是$C^\infty$的，
即Jacobi矩阵$\frac{\partial {x}^j_\beta}{\partial {x}^i_\alpha}$是$\{x^i_\alpha\}$的$C^\infty$函数，
同样，其逆矩阵$\frac{\partial {x}^i_\alpha}{\partial {x}^j_\beta}$也是$\{{x}^j_\beta\}$的$C^\infty$函数。
设$p \in U_\alpha \cap U_\beta$，则式\eqref{chfb:eqn_ctb_atalas}中的$\omega$在$p$点的
表达式有两个等价描述如下：
\begin{equation*}
    \omega = \omega_i \left. \d x^i_\alpha \right | _p
    = \tilde{\omega}_j \left. \d x^j_\beta \right | _p
\end{equation*}
不难得到两者分量的变换关系：
\begin{equation}\label{chfb:eqn_ctbv}
    \tilde{\omega}_j = \omega_i \left. \frac{\partial {x}^i_\alpha}{\partial {x}^j_\beta}  \right | _p
    {\quad \Leftrightarrow \quad}
    \omega_i  = \tilde{\omega}_j \left. \frac{\partial {x}^j_\beta}{\partial {x}^i_\alpha}  \right | _p .
\end{equation}
上述坐标变换关系中的Jacobi矩阵是$C^\infty$的，坐标$(\omega_1,\cdots,\omega_m)$本身也是$C^\infty$的，
那么坐标变换$\omega_i \leftrightarrow \tilde{\omega}_j$是$C^\infty$的；
而$U_\alpha \cap U_\beta$自身的坐标变换$x^i_\alpha \leftrightarrow {x}^j_\beta$本来就是$C^\infty$的；
$T^*M$坐标是由两者结合得到的$(x^1,\cdots,x^m,\omega_1,\cdots,\omega_m)$，它的坐标变换自然也是$C^\infty$的。
因此这些坐标卡构成光滑图册，将所有与之相容的坐标卡加入便得到唯一的最大光滑图册，
从而建立了$T^*M$的$C^\infty$微分结构（见定义\ref{def:smooth-manifold}）。

由式\eqref{chfb:eqn_ctbv}可以直接看出余切丛的过渡函数族是
\begin{equation}
    g_{\alpha\beta} = \frac{\partial {x}^j_\beta}{\partial {x}^i_\alpha}  .
    \qquad \text{注意}g_{\alpha\beta}\text{是} m\times m \text{阶矩阵}
\end{equation}
上式恰好是切丛过渡函数族的转置逆。
\qed

\begin{example}
    矢量丛$E$的对偶丛$E^*$。
\end{example}
设$E\to M$是以$V$为纤维型的矢量丛。它的\textbf{对偶丛}逐纤维定义为
$E^*:=\bigsqcup_{p\in M}E_p^*$，丛投影记为$\tilde\pi:E^*\to M$；其中$E_p:=\pi^{-1}(p)$。
设$\psi_{\alpha,p}:E_p\to V$是局部平凡化在纤维上的限制，
对偶平凡化$\tilde\psi_{\alpha,p}:E_p^*\to V^*$由
$\tilde\psi_{\alpha,p}(\lambda):=\lambda\circ\psi_{\alpha,p}^{-1}$给出。
在两个平凡化的交集上，若
\begin{equation}
    \psi_{\alpha,p}^{-1}(y_\alpha)=\psi_{\beta,p}^{-1}(y_\beta),\qquad
    \tilde\psi_{\alpha,p}^{-1}(\lambda_\alpha)=\tilde\psi_{\beta,p}^{-1}(\lambda_\beta),
\end{equation}
则同一配对的数值满足
\begin{equation}\label{chfb:eqn_ylaylb}
    \lambda_\alpha(y_\alpha)=\lambda_\beta(y_\beta).
\end{equation}
也就是说，纤维间的自然配对为
\begin{equation}\label{chfb:eqn_dph}
    \bigl(\tilde\psi_{\alpha,p}^{-1}(\lambda_\alpha)\bigr)
    \bigl(\psi_{\alpha,p}^{-1}(y_\alpha)\bigr):=\lambda_\alpha(y_\alpha).
\end{equation}
取$V$的一组基及$V^*$的对偶基，将$y=y^i e_i$的分量记成列，
将$\lambda=\lambda_i e^i$的分量记成行，则
\[
    \lambda(y)=\lambda_i y^i=(\lambda_1,\ldots,\lambda_q)(y^1,\ldots,y^q)^{\mathrm{T}}.
\]
由$y_\beta=g_{\beta\alpha}^p y_\alpha$及式\eqref{chfb:eqn_ylaylb}可知，
对偶矢量的行分量满足
\[
    \lambda_\beta=\lambda_\alpha(g_{\beta\alpha}^p)^{-1}
    =\lambda_\alpha g_{\alpha\beta}^p.
\]
若把对偶矢量的分量也排列成列，则对偶丛的过渡矩阵为
\[
    h_{\beta\alpha}:=(g_{\alpha\beta}^p)^{\mathrm{T}}=(g_{\beta\alpha}^p)^{-\mathrm{T}}.
\]
这些矩阵满足过渡函数的相容条件，因而确定$E^*$的光滑矢量丛结构。
式\eqref{chfb:eqn_dph}与平凡化的选择无关，因为用行分量计算有
\[
    \lambda_\beta(y_\beta)
    =(\lambda_\alpha g_{\alpha\beta}^p)(g_{\beta\alpha}^p y_\alpha)
    =\lambda_\alpha(y_\alpha).
\]
该配对逐纤维双线性且非退化。\qed

\begin{example}\label{chfb:exm_oplus}
    矢量丛的直和。
\end{example}
设$\pi: E \rightarrow M$、$\tilde{\pi}: \widetilde{E} \rightarrow M$是$M$上的两个矢量丛，
纤维型空间分别是$V$、$\tilde{V}$，
过渡函数族分别是$\left\{g_{\alpha \beta}\right\}$、$\left\{\tilde{g}_{\alpha \beta}\right\}$。 命
\begin{equation*}
    h_{\alpha \beta}=\begin{pmatrix}
        g_{\alpha \beta} & 0 \\ 0 & \tilde{g}_{\alpha \beta}
    \end{pmatrix}.
\end{equation*}
它是左作用在直和空间$V \oplus \tilde{V}$上的线性自同构，对于 $y \in V$、$ \tilde{y} \in \tilde{V}$，有
\begin{equation*}
    h_{\alpha \beta} \cdot (y_\beta \oplus \tilde{y}_\beta)
    =h_{\alpha \beta} \cdot \begin{pmatrix}
        y_{\beta} & 0 \\ 0 & \tilde{y}_{\beta}
    \end{pmatrix}
    =\begin{pmatrix}
        g_{\alpha \beta}\cdot  y_{\beta} & 0 \\ 0 & \tilde{g}_{\alpha \beta}\cdot  \tilde{y}_{\beta}
    \end{pmatrix}
    =\begin{pmatrix}
        y_{\alpha} & 0 \\ 0 & \tilde{y}_{\alpha}
    \end{pmatrix} .
\end{equation*}
函数族$\{h_{\alpha \beta}: U_\alpha \cap U_\beta \rightarrow GL(V \oplus \tilde{V})\}$满足
命题\ref{chfb:thm_trans-func}中的条件(1)和(2)；
因此，根据定理\ref{chfb:thm_vector-existence}在$M$上存在以$\left\{h_{\alpha \beta}\right\}$为过渡函数族、
以$V\oplus \tilde{V}$为纤维型的矢量丛，记作$E\oplus\widetilde{E}$，称为矢量丛$E$和$\widetilde{E}$的直和。
\qed

\begin{example}\label{chfb:exm_otimes}
    矢量丛的张量积。
\end{example}
假定矢量丛$E$、$\tilde{E}$如例\ref{chfb:exm_oplus}所述。
用$h_{\alpha \beta}$表示线性变换$g_{\alpha \beta}$和$\tilde{g}_{\alpha \beta}$的张量积，
它左作用在张量积$V \otimes \tilde{V}$上，其定义是：对于 $v \in V$、$ \tilde{v} \in \tilde{V}$，有
\begin{equation*}
    h_{\alpha \beta} \cdot(v \otimes \tilde{v})
    :=\left(g_{\alpha \beta} \cdot v\right) \otimes
    \left(\tilde{g}_{\alpha \beta} \cdot \tilde{v}\right) .
\end{equation*}
容易验证，函数族$\{h_{\alpha \beta}: U_\alpha \cap U_\beta \rightarrow GL(V \otimes \tilde{V})\}$满足
命题\ref{chfb:thm_trans-func}中的条件(1)、(2)。
因此，根据定理\ref{chfb:thm_vector-existence}在$M$上存在以$\left\{h_{\alpha \beta}\right\}$为过渡函数族、
以$V\otimes\tilde{V}$为纤维型的矢量丛，记作$E \otimes \tilde{E}$，称为矢量丛$E$和$\tilde{E}$ 的张量积。
\qed

\begin{example}
    光滑流形上的张量丛。
\end{example}
在例\ref{chfb:exm_tb}、\ref{chfb:exm_ctb}已经给出切丛、余切丛的示例，
下面用矢量丛的张量积来构造一般的张量丛；这也是例\ref{chfb:exm_otimes}的一个应用。
我们仅以$(1,1)$型张量丛为例来说明。
令$V^1_1$表示$\mathbb{R}^m$上的$(1,1)$型张量空间；
显然它可以表示成
\begin{equation*}
    V^1_1 = \operatorname{span}_{\mathbb{R}}\left\{  \frac{\partial}{\partial x^k}\otimes \d x^j \right\} .
\end{equation*}
这是切空间和余切空间的张量积。过渡函数族是
\begin{equation*}
    g_{\alpha\beta} = \frac{\partial {x}^l_\alpha}{\partial {x}^k_\beta} \otimes
    \frac{\partial {x}^j_\beta}{\partial {x}^i_\alpha}  .
    \qquad \text{矩阵张量积；}\alpha,\beta \text{不求和}
\end{equation*}
有了纤维型和过渡函数族，由例\ref{chfb:exm_otimes}可知
存在由$1$个切丛与$1$个余切丛构成的张量积，它就是$(1,1)$型张量丛。
这个张量丛自然可以拓展到$(r,s)$型，即
\begin{equation*}
    T_s^r(M)=\bigcup_{p \in M} T_s^r(p); \qquad
    T_s^r(p)\text{是}p\in M\text{处的}(r,s)\text{型张量空间}
\end{equation*}

张量丛的光滑截面是光滑张量场。
\qed

\section{矢量丛的联络与曲率}

矢量丛联络给出截面的协变导数。取 $E=TM$ 时，它就是第\ref{appx:curvature}章仿射联络在切矢量场上的限制。

\begin{definition}\label{chfb:def_vb-conncection}
    设$(E,M,\pi)$是光滑流形$M$上的一个矢量丛，$\Gamma(E)$是该矢量丛光滑截面的集合。
    矢量丛 $E$上的一个 \textbf{联络} $\nabla$是指存在映射
    $\nabla: \mathfrak{X}(M) \times \Gamma(E) \rightarrow \Gamma(E)$。
    并且此映射满足以下条件：$\forall X, Y \in \mathfrak{X}(M)$，
    $\forall \xi, \eta \in \Gamma(E)$，$\forall c \in \mathbb{R}$，
    及 $\forall f \in C^{\infty}(M)$，有如下三条公理成立
    （令$\nabla(X, \xi) := \nabla_X \xi$ ，\textbf{协变导数}）。

    (1) $\mathbb{R}$线性： $\nabla_X(\xi+c \eta) =
     \nabla_X \xi+\nabla_X (c\eta) =  \nabla_X \xi+c \nabla_X \eta$；

    (2) 关于方向的 $C^\infty(M)$-线性： $\nabla_{X+f Y} \xi =
    \nabla_X \xi+ \nabla_{fY} \xi = \nabla_X \xi+f \nabla_Y \xi$；

    (3) $\nabla_X(f \xi) = (\nabla_X f) \xi+f \nabla_X \xi
     = X(f) \xi+f \nabla_X \xi$。
\end{definition}

定义\ref{chfb:def_vb-conncection}中的条件(2)表明 $\nabla_X \xi$ 关于自变量 $X$ 有张量性质，
于是可以把联络$\nabla$ 看作映射 $\nabla: \Gamma(E) \rightarrow \Gamma (T^* M \otimes E )$，使得
$(\nabla \xi)(X)=\nabla_X \xi$，$ \forall \xi \in \Gamma(E), X \in \mathfrak{X}(M) $成立。

矢量丛联络总是存在。

\begin{theorem}
    设$(E,M,\pi)$是任意矢量丛，则矢量丛$E$上必定存在一个联络。
\end{theorem}
\begin{proof}
    取局部平凡化开覆盖$\{U_\alpha\}$及其从属的局部有限单位分解$\{\rho_\alpha\}$。
    在$U_\alpha$的局部标架中定义$\nabla^{(\alpha)}$为逐分量微分，并令
    $\nabla\xi=\sum_\alpha\rho_\alpha\nabla^{(\alpha)}\xi$；各项在$U_\alpha$外延拓为零。
    局部有限性保证该式光滑；线性性由各局部联络继承，而
    \[
        \nabla(f\xi)=\sum_\alpha\rho_\alpha
        (\d f\otimes\xi+f\nabla^{(\alpha)}\xi)
        =\d f\otimes\xi+f\nabla\xi .
    \]
    因而$\nabla$满足联络的定义。
\end{proof}

矢量丛联络也具有局部性。
\begin{proposition}
    设 $\nabla$ 是矢量丛 $(E,M,\pi)$ 上的一个联络，$\xi, \eta \in \Gamma(E)$，
    $ X, Y \in \mathfrak{X}(M)$，$U$ 是 $M$上的一个开集。
    若$ \xi|_U = \eta|_U$，$ X|_U = Y|_U $；
    则 $\left.\left(\nabla_X \xi\right)\right|_U=\left.\left(\nabla_Y \eta\right)\right|_U$。
\end{proposition}

\subsection{张量丛联络}

切丛联络通过对偶配对和张量积的 Leibniz 法则，唯一延拓为各阶张量丛上的联络，
与第\ref{appx:curvature}章的仿射联络定义一致。
下面给出直和、张量积及对偶丛上的诱导联络。
设 $\nabla^{(i)}$ 是矢量丛 $(E_i, M,\pi_i)$（$i=1,2$） 上的联络。
由 $E_1$ 和 $E_2$ 可以构造矢量丛的直和 $E_1 \oplus E_2$ 和张量积 $E_1 \otimes E_2$。
我们要给出 $\nabla^{(1)}$ 和 $\nabla^{(2)}$ 在这两个矢量丛上的诱导联络；
对于任意的 $X \in \mathfrak{X}(M)$，和对于任意的 $\xi_i \in \Gamma(E_i)$，
定义：
\begin{align}
    \nabla_X\left(\xi_1 \oplus \xi_2\right):=&
    \left(\nabla_X^{(1)} \xi_1\right) \oplus \nabla_X^{(2)} \xi_2,  \\
    \nabla_X\left(\xi_1 \otimes \xi_2\right):=& \left(\nabla_X^{(1)} \xi_1\right) \otimes
    \xi_2+\xi_1 \otimes \nabla_X^{(2)} \xi_2 . \label{chfb:eqn_otcon}
\end{align}
容易验证：上面两个映射满足定义\ref{chfb:def_vb-conncection}中的三个公理，
因而是矢量丛 $E_1 \oplus E_2$ 和 $E_1 \otimes E_2$ 上的联络。
我们规定式\eqref{chfb:eqn_otcon}中$\nabla_X$与缩并运算可对易。

\begin{example}
    对偶丛$(E^*,M,\tilde{\pi})$上的联络$\nabla^*$（导出后，将$\nabla^*$简记为$\nabla$）。
\end{example}
设$\alpha \in \Gamma(E)$，$\alpha^*\in \Gamma(E^*)$，
两者的自然配对为$\alpha^*(\alpha)$。
目前我们并不清楚$\nabla_X^* \alpha^*$是如何定义的，下面导出它。
对于任意的 $X \in \mathfrak{X}(M)$，有
\begin{equation*}
    \nabla_X( \alpha^* \otimes \alpha )\xlongequal{\ref{chfb:eqn_otcon}}
    (\nabla_X^* \alpha^*)\otimes \alpha + \alpha^* \otimes \nabla_X \alpha .
\end{equation*}
令$\alpha^*$与$\alpha$缩并（也就是求配对），
可得标量函数场$C(\alpha^* \otimes \alpha)=\alpha^*(\alpha)$；
而$\nabla_X\bigl( \alpha^*(\alpha) \bigr)=X \bigl( \alpha^*(\alpha) \bigr)$；
再利用缩并与$\nabla_X$可对易，由上式得
\begin{equation}\label{chfb:eqn_dual-con}
    (\nabla_X^* \alpha^* )(\alpha)=X\bigl(\alpha^*(\alpha)\bigr)-\alpha^*\bigl(\nabla_X \alpha\bigr).
\end{equation}
通过上面操作得到了一个映射$(\alpha^*, X) \mapsto \nabla_X^* \alpha^*$，
容易看出映射$\nabla^*$是从笛卡尔积$\mathfrak{X}(M)\times \Gamma(E^*)$到
对偶丛截面$\Gamma (E^*)$的一个联络（请读者验证式\eqref{chfb:eqn_dual-con}中的$\nabla_X^*$满足
矢量丛联络定义中的三个公理）；称$\nabla^*$为与矢量丛联络 $\nabla$相互对偶的对偶丛联络。
\qed

有了张量积上联络定义\eqref{chfb:eqn_otcon}以及对偶丛上的联络\eqref{chfb:eqn_dual-con}，
我们再来看看张量丛上联络与切丛联络的关系。
注意到切丛、余切丛相互对偶，故从联络\eqref{chfb:eqn_dual-con}可得余切丛的联络。
现在有了切矢量场的联络，以及余切矢量场的联络，我们利用式\eqref{chfb:eqn_otcon}可以
得到任意张量丛 $T_q^p(M)$上的联络。
张量丛$T_q^p(M)$是$p$个$TM$与$q$个$T^* M$构成的张量积矢量丛，
对于任意的 $\tau \in \Gamma (T_q^p(M)) =\mathfrak{T}_q^p(M)$，
其作用可直接写为
\begin{align*}
    &(\nabla_X\tau)(\omega_1,\ldots,\omega_p;v_1,\ldots,v_q)\\
    &\quad=X\bigl(\tau(\omega_1,\ldots,\omega_p;v_1,\ldots,v_q)\bigr)\\
    &\qquad-\sum_{r=1}^{p}\tau(\omega_1,\ldots,\nabla_X\omega_r,\ldots,\omega_p;v_1,\ldots,v_q)\\
    &\qquad-\sum_{r=1}^{q}\tau(\omega_1,\ldots,\omega_p;v_1,\ldots,\nabla_Xv_r,\ldots,v_q).
\end{align*}
在局部标架下展开，即得通常的张量协变导数分量公式。

因此，切丛联络与前文作用于所有张量场的仿射联络相互唯一确定。

\subsection{局部标架表示}\label{chfb:sec_vc}

设$(E,M,\pi)$是一个秩为$q$的矢量丛，$U$是$M$的任意一个非空开子集，
$\tilde{\pi}$是$\pi$在 $\pi^{-1}(U)$上的限制，
则有矢量丛$E$在$U$上的限制：$\tilde{\pi}:\left.E\right|_U=\pi^{-1}(U) \to U$。

利用矢量丛联络的局部性定理不难知道，从$E$上的联络$\nabla$可以诱导出局部开集$\left.E\right|_U$上的联络。
若在$U\subset M$上存在局部切标架场$\left\{e_i,\ 1 \leqslant i \leqslant m\right\}$，
其余切标架场是$\left\{e^{i},\ 1 \leqslant i \leqslant m\right\}$；
以及矢量丛$\left.E\right|_U$上的局部标架场$\left\{s_\alpha,\ 1 \leqslant \alpha \leqslant q\right\}$，
$s^{\alpha}$是其对偶标架场；那么对于 $X \in \mathfrak{X}(M)$，$ \xi \in \Gamma(E)$，可设
\begin{equation*}
    \left.X\right|_U= \sum_{i=1}^{m}X^i e_i,\left.\quad
    \xi\right|_U= \sum_{\alpha=1}^{q}\xi^\alpha s_\alpha,
    \qquad \text{其中} \  X^i, \xi^\alpha \in C^{\infty}(U) .
\end{equation*}
沿用第\ref{appx:curvature}章的约定：联络系数的第一个下标表示求导方向，第二个下标表示被求导的基矢。定义
\begin{equation}\label{chfb:eqn_con-Gamma}
    \nabla_{e_i}s_\alpha=\Gamma^\beta_{i \alpha} s_\beta,\qquad
    \Gamma^\beta_{i \alpha}:=s^{\beta}(\nabla_{e_i}s_\alpha)\in C^\infty(U).
\end{equation}
$\Gamma^\beta_{i \alpha}$称为联络$\nabla$在标架场
$\left\{e_i\right\}$及$\left\{s_\alpha\right\}$下的\textbf{联络系数}。

根据联络的条件，有
\begin{equation}\label{chfb:eqn_con-comp}
\begin{aligned}
    \left. (\nabla_X \xi )\right|_U =& \nabla_{X|_U} (\xi |_U ) = X^i \nabla_{e_i}
    (\xi^\alpha s_\alpha ) = X^i \bigl((\nabla_{e_i} \xi^\alpha ) s_\alpha+
    \xi^\alpha \nabla_{e_i} s_\alpha \bigr) \\
    \xlongequal{\ref{chfb:eqn_con-Gamma}}&
    \sum_{\alpha=1}^{q}\sum_{i=1}^{m}
    X^i\left(e_i(\xi^\alpha)+\sum_{\beta=1}^{q}\xi^\beta\Gamma^\alpha_{i \beta}\right)s_\alpha .
\end{aligned}
\end{equation}
式\eqref{chfb:eqn_con-comp}给出了协变导数的局部表达式。定义联络1-形式
\begin{equation}\label{chfb:eqn_con-form}
    \omega^\alpha_{\beta}:=\Gamma^\alpha_{i \beta}e^{i},\qquad
    \omega^\alpha_{\beta}(e_i)=\Gamma^\alpha_{i \beta}.
\end{equation}
对每一对固定的$\alpha,\beta$，$\omega^\alpha_{\beta}$都是底流形$U$上的1-形式，
其自变量属于$TU$；$\alpha,\beta$则是丛$E|_U$的标架指标。
协变微分$\nabla\xi$是取值于$E$的1-形式，视为$T^*M\otimes E$的截面，约定
$(\nabla\xi)(X):=\nabla_X\xi$。于是
\begin{equation}\label{chfb:eqn_Dsos}
    \nabla s_\alpha=\omega^\beta_{\alpha}\otimes s_\beta,\qquad
    \nabla_{e_i}s_\alpha=\omega^\beta_{\alpha}(e_i)s_\beta.
\end{equation}
一般地，
\begin{equation}
    \left.(\nabla\xi)\right|_U
    =\bigl(\d \xi^\beta+\xi^\alpha\omega^\beta_{\alpha}\bigr)\otimes s_\beta.
\end{equation}
这里$(\eta\otimes s)(X)=\eta(X)s$，从而上述公式与式\eqref{chfb:eqn_con-comp}一致。

若$\left\{\tilde{s}_\alpha\right\}$ 是矢量丛 $E$ 在 $U$ 上的另一个局部标架场，设
\begin{equation}\label{chfb:eqn_dsom}
    \tilde{s}_\alpha=a_{\alpha}^\beta s_\beta, \qquad
    \nabla \tilde{s}_\alpha=\tilde{\omega}_{\alpha}^\beta\otimes\tilde{s}_\beta .
    \quad \text{其中}\ a_{\alpha}^\beta \in C^{\infty}(U)
\end{equation}
在$U$上的每一点处 $\operatorname{det}\left(a_{\alpha}^\beta\right) > 0$，
用以保证标架场定向相同。则有
\begin{equation}
    \nabla \tilde{s}_\alpha=\d a_{\alpha}^\beta\otimes s_\beta
    +a_{\alpha}^\beta \nabla s_\beta
    =\left(\d a_{\alpha}^\beta
    +a_{\alpha}^\gamma \omega_{\gamma}^\beta\right)\otimes s_\beta .
\end{equation}
与式\eqref{chfb:eqn_dsom}的第二式相比较，可得
\begin{equation}\label{chfb:eqn_wform-trans}
	a_{\gamma}^\beta \tilde{\omega}_{\alpha}^\gamma
	=\d a_{\alpha}^\beta +
	\omega_{\gamma}^\beta a_{\alpha}^\gamma
	\ \Leftrightarrow \
	\tilde{\omega}_{\alpha}^\sigma=(a^{-1})_{\beta}^{\sigma}
	\d a_{\alpha}^\beta+(a^{-1})_{\beta}^{\sigma}
	\omega_{\gamma}^\beta a_{\alpha}^\gamma .
\end{equation}
式\eqref{chfb:eqn_wform-trans}是联络形式在局部标架场变换下的变换公式。
此式在切丛上对应于定理\ref{thm:connection-frame-change}中的仿射联络换基公式；
式\eqref{chfb:eqn_wform-trans}适用于矢量丛（包含切标架丛）。

\begin{example}[沿曲线的平行移动]\label{chfb:eg_parallel-lift}
    设$\nabla$为矢量丛$\pi:E\to M$上的联络，$\gamma:[a,b]\to M$为光滑曲线。
    沿线矢量$v(t)\in E_{\gamma(t)}$就是$\gamma$在$E$中的提升；它只须沿参数$t$定义，不必延拓为$M$上的截面。
    在局部标架中写$v(t)=v^\alpha(t)s_\alpha(\gamma(t))$，则
    \eq{\frac{D v}{\d t}
        =\left(\dot v^\beta+\omega^\beta_{\alpha}(\dot\gamma)v^\alpha\right)s_\beta.}
    此式由联络的换基律保证与标架选择无关。若$D v/\d t=0$，就称$v$沿$\gamma$\textbf{平行}，其分量满足
    \eq{\label{chfb:eq_parallel-lift}
        \dot v^\beta+\omega^\beta_{\alpha}(\dot\gamma)v^\alpha=0.}
    这是线性常微分方程组，给定$v(a)$后，在整个$[a,b]$上有唯一解；跨越标架域时按过渡函数拼接。
    因而\textbf{平行移动}映射
    \eq{P_\gamma:E_{\gamma(a)}\to E_{\gamma(b)},\qquad v(a)\mapsto v(b)}
    是线性同构，其逆由反向曲线给出。它一般依赖整条路径，而非仅依赖端点。
    从总空间看，平行移动正是沿特定提升将初始纤维中的点送到终止纤维；这种提升将在定义\ref{chfb:def_horizontal-lift}中称为水平提升。
\end{example}

\begin{example}
    对偶丛联络的局部表示。
\end{example}
把式\eqref{chfb:eqn_dual-con}的截面分别取成标架，$X\to {e_i}$，有
\begin{equation}
    (\nabla_{e_i} s^{\alpha} )(s_\beta)= {e_i} \bigl(s^{\alpha}(s_\beta)\bigr)
    -s^{\alpha} \bigl(\nabla_{e_i} s_\beta\bigr)
    =-s^{\alpha} \bigl( \Gamma^\gamma_{i \beta} s_\gamma \bigr)
    =-\Gamma^\alpha_{i \beta} .
\end{equation}
在切丛上，这就是对偶基的协变导数公式 $(\nabla_{e_i}e^\alpha)(e_\beta)=-\Gamma^\alpha_{i\beta}$，可由 $e^\alpha(e_\beta)=\delta^\alpha_\beta$ 的协变导数为零得到。
\qed

\subsection{拉回丛上的诱导联络}

设$M$、$N$是光滑流形，维数分别为$m$、$n$；$\phi:M\to N$是光滑浸入映射。命
\begin{equation}
    \phi^* TN := \bigcup_{p \in M} \{p\}\times T_{\phi(p)} N .
\end{equation}
则$\phi^* TN$是流形$M$上秩为$n$的矢量丛（证明见后）；称为切丛$TN$通过映射$\phi$拉回
到$M$上的矢量丛，简称\textbf{拉回丛}。
拉回丛$\phi^* TN$的截面称为$N$中\textbf{沿映射$\phi$定义的切矢量场}。

下面证明$\phi^* TN$是流形$M$上秩为$n$的矢量丛。
定义丛投影$\pi : \phi^* TN \to M$使得对于任意$p\in M$
有$\pi\bigl(\{p\}\times T_{\phi(p)} \bigr)=\{p\} $。
任取$N$的一族局部坐标系$\mathcal{V}=\{(V_a; y^\lambda)\mid a \in J\}$，
使得$\phi(M)\subset \cup_{a\in J} V_a$。
对每个$p\in M$，取$N$中含$\phi(p)$的坐标邻域$(V_a,y_a^\lambda)$，
再取$M$的坐标邻域$(U_\alpha,x_\alpha^i)$，使$\phi(U_\alpha)\subset V_a$。
定义局部平凡化
\[
 \psi_\alpha\left(q,v^\lambda
 \left.\frac{\partial}{\partial y_a^\lambda}\right|_{\phi(q)}\right)
 =(q,(v^1,\ldots,v^n)),\qquad q\in U_\alpha .
\]
这些映射的定义域$\pi^{-1}(U_\alpha)$覆盖$\phi^*TN$。
在两组坐标的交叠处，纤维坐标满足
\[
 v_b^\mu=\frac{\partial y_b^\mu}{\partial y_a^\lambda}(\phi(q))v_a^\lambda .
\]
此变换光滑且对纤维坐标线性，逆变换由逆坐标变换的Jacobi矩阵给出。
因此这些局部平凡化赋予$\phi^*TN$秩为$n$的光滑矢量丛结构。

下面讨论{拉回丛上的诱导联络}。设$N$上有仿射联络$\mathrm{D}$，
借助它能够在$\phi^* T N$上引入诱导联络$\nabla$，具体构造如下。
对于任意$\xi \in \Gamma(\phi^* T N)$ 和 $p \in M$，取 $N$ 在点 $\phi(p) \in N$的
开邻域 $V$ 以及 $M$ 在点 $p$ 的邻域 $U$ 使得 $\phi(U) \subset V$，并在 $V$ 上取局部标架
场 $\left\{e_\alpha\right\}$。则在邻域 $U$ 上 $\xi$ 可以表示为
$\xi |_U=\xi^\alpha \cdot e_\alpha$，
其中$\xi^\alpha \in C^{\infty}(U)(1 \leqslant \alpha \leqslant n)$。
对于任意的 $X \in T_p M$，构造\textbf{拉回丛上的诱导联络}为：
\begin{equation}\label{chfb:eqn_pbb-connection}
    \nabla_X \xi := X(\xi^\alpha ) e_\alpha\bigl(\phi(p)\bigr)+\xi^\alpha(p)
    \mathrm{D}_{\phi_*(X(p))} e_\alpha\bigl(\phi(p)\bigr), \quad \in \pi^{-1}(p) .
\end{equation}
下文将验证上式右端与局部标架场 $\left\{e_\alpha\right\}$ 的选取无关。
如果 $X \in \mathfrak{X}(M)$，则式\eqref{chfb:eqn_pbb-connection}给出的$\nabla_X \xi$是
拉回丛 $\pi: \phi^* T N \rightarrow M$ 的光滑截面。
于是，得到映射 $\nabla: \mathfrak{X}(M) \times \Gamma\left(\phi^* T N\right)
\rightarrow \Gamma\left(\phi^* T N\right)$。
不难验证，式\eqref{chfb:eqn_pbb-connection}给出的映射$\nabla$满足
定义\ref{chfb:def_vb-conncection}中的所有条件，
因而是矢量丛 $\phi^* T N$ 上的一个联络。

现证明此定义与局部标架场$\{e_\alpha\}$选取无关。
设另有标架场$\{\tilde{e}_{\sigma}\}$，它们之间存在变换关系
$\tilde{e}_\sigma= e_\alpha A^\alpha_{\sigma} {\  \Leftrightarrow\  }
\tilde{e}_\sigma B_{\alpha}^\sigma  = e_\alpha$（$A=B^{-1}$）。
则$\phi$上矢量场$\xi$可展为
\begin{equation}
    \xi = e_\alpha \xi^\alpha = \tilde{e}_\sigma \tilde{\xi}^\sigma
    = e_\alpha A^\alpha_{\sigma} \tilde{\xi}^\sigma
    = \tilde{e}_\sigma B_{\alpha}^\sigma \xi^\alpha .
\end{equation}
将上式代入诱导联络\eqref{chfb:eqn_pbb-connection}，
有（为节省空间，省略后缀$\phi(p)$、$p$等；拉回丛$\phi^*TN$上
的标架场$e_\alpha$和变换系数$A$、$B$在$\phi(p)$处取值，系数$\xi^\alpha$在$p$点取值）
\begin{align*}
    \nabla_X \xi =& X(\xi^\alpha ) e_\alpha\bigl(\phi(p)\bigr)
    +\xi^\alpha(p) \mathrm{D}_{\phi_*(X(p))} e_\alpha\bigl(\phi(p)\bigr) \\
    =& X\left(A^\alpha_{\sigma} \tilde{\xi}^\sigma \right)
    \tilde{e}_\rho B_{\alpha}^\rho
    + A^\alpha_{\sigma} \tilde{\xi}^\sigma
    \mathrm{D}_{\phi_*X} \bigl(\tilde{e}_\rho B_{\alpha}^\rho\bigr) \\
    =& X\left(A^\alpha_{\sigma}  \right) \tilde{\xi}^\sigma
    \tilde{e}_\rho B_{\alpha}^\rho
    + X ( \tilde{\xi}^\rho ) \tilde{e}_\rho
    +  \tilde{\xi}^\rho  \mathrm{D}_{\phi_*X} \bigl(\tilde{e}_\rho \bigr)
    + A^\alpha_{\sigma} \tilde{\xi}^\sigma \tilde{e}_\rho
    X( B_{\alpha}^\rho \circ \phi ) \\
    =& \tilde{\xi}^\sigma \tilde{e}_\rho \bigl(
    X(A^\alpha_{\sigma}  ) B_{\alpha}^\rho
    + A^\alpha_{\sigma} X( B_{\alpha}^\rho ) \bigr)
    + X ( \tilde{\xi}^\rho ) \tilde{e}_\rho
    +  \tilde{\xi}^\rho  \mathrm{D}_{\phi_*X} \bigl(\tilde{e}_\rho \bigr)   \\
    =& \tilde{\xi}^\sigma \tilde{e}_\rho  X( \delta^\rho_\sigma )+
    X ( \tilde{\xi}^\rho ) \tilde{e}_\rho
    +  \tilde{\xi}^\rho  \mathrm{D}_{\phi_*X} \bigl(\tilde{e}_\rho \bigr)
    =X ( \tilde{\xi}^\rho ) \tilde{e}_\rho
    +  \tilde{\xi}^\rho  \mathrm{D}_{\phi_*X} \bigl(\tilde{e}_\rho \bigr) .
\end{align*}
这便证明了诱导联络定义与标架场选取无关。

注意到，对于任意的 $X \in \mathfrak{X}(M), \phi_*(X) \in \Gamma(\phi^* T N)$。
现在进一步假定 $N$是广义黎曼流形，则$\mathrm{D}$是相应的Levi-Civita联络（见定义\ref{def:levi-civita}）。
则根据联络 $\nabla$ 的构造以及Levi-Civita联络$\mathrm{D}$的性质，
可得下面两个恒等式：
\begin{align}
    X\langle\xi, \eta\rangle= &\left\langle\nabla_X \xi, \eta\right\rangle+\left\langle\xi,
     \nabla_X \eta\right\rangle, \ \forall X \in \mathfrak{X}(M),\ \xi, \eta \in \Gamma(\phi^* T N)
      \\
    \phi_*([X, Y])=&\nabla_X \phi_*(Y)-\nabla_Y \phi_*(X), \quad \forall X, Y \in \mathfrak{X}(M) .
\end{align}
第一个式子说明度规与诱导联络是相容的，第二个式子说明诱导联络是无挠的。
\begin{proof}
取$N$的局部标架，将$\xi=\xi^\alpha e_\alpha$、$\eta=\eta^\beta e_\beta$代入。
对$\langle\xi,\eta\rangle=\xi^\alpha\eta^\beta(g_{\alpha\beta}\circ\phi)$求导，
使用$\mathrm D g=0$，乘积法则中的三项正好合为
$\langle\nabla_X\xi,\eta\rangle+\langle\xi,\nabla_X\eta\rangle$。
对于第二式，记$\phi^\mu=y^\mu\circ\phi$，在坐标标架中有
\[
 \bigl(\nabla_X\phi_*Y\bigr)^\lambda
 =X(Y\phi^\lambda)
  +(\Gamma^\lambda_{\mu\nu}\circ\phi)(X\phi^\mu)(Y\phi^\nu).
\]
交换$X,Y$相减，因$\Gamma^\lambda_{\mu\nu}=\Gamma^\lambda_{\nu\mu}$，
联络项相消，剩余$[X,Y]\phi^\lambda$即$\phi_*[X,Y]$的分量。
\end{proof}

诱导联络还可用局部延拓理解。因$\phi$是浸入，可缩小$U\subset M$使$\phi|_U$为嵌入，
将沿$\phi$的截面$\xi$延拓为$N$上的局部矢量场$\bar\xi$，则
\begin{equation}
    (\nabla_X\xi)_p=\left.\mathrm D_{\phi_{*p}X_p}\bar\xi\right|_{\phi(p)} .
\end{equation}
若两个延拓沿$\phi(U)$相同，其差的分量在$\phi(U)$上为零，沿切于$\phi(U)$的方向求导也为零，
故右端与延拓的选择无关。这正是沿浸入曲线求协变导数的含义。

\begin{example}
    浸入曲线例子。
\end{example}
设$\gamma:[a, b] \rightarrow N$是仿射联络空间 $(N, \mathrm{D})$上的正则光滑曲线，
即切矢量 $\gamma^{\prime}$ 处处不为零；则 $\gamma$ 在局部上是正则嵌入。
因此，对于任意的 $X \in \Gamma(\gamma^* TN)$，$X$在局部上是$N$上的
切矢量在曲线 $\gamma$ 上的限制，因而协变导数 $\mathrm{D}_{\gamma^{\prime}(t)} X$ 处处有意义。
根据诱导联络 $\nabla$ 和 $\mathrm{D}_{\gamma^{\prime}(t)} X$的定义可知
\begin{equation}
    \nabla_{\frac{\partial}{\partial t}} X=\mathrm{D}_{\gamma^{\prime}(t)} X,
    \qquad \forall X \in \Gamma(\gamma^* T N) .
\end{equation}
据此，$X$ 沿曲线 $\gamma$ 是平行矢量场的条件可以改写为$\nabla_{\frac{\partial}{\partial t}} X=0$。
特别地，$\gamma$是测地线等价于$\nabla_{\frac{\partial}{\partial t}} \gamma^{\prime}=0$。
\qed

\subsection{曲率}

将式\eqref{chfb:eqn_wform-trans}改成矩阵形式；
令$\omega=(\omega_{\gamma}^\beta)$、
$\tilde{\omega}=(\tilde{\omega}_{\gamma}^\beta)$、
$A=(a_{\gamma}^\beta)$，则有
\begin{equation}\label{chfb:eqn_wftmat}
    A \tilde{\omega}= \d A + \omega A
    \quad \Leftrightarrow \quad
    \tilde{\omega}= A^{-1} \d A + A^{-1}\omega A .
\end{equation}
式\eqref{chfb:eqn_wftmat}比式\eqref{chfb:eqn_wform-trans}简洁了许多。
下面，我们将利用式\eqref{chfb:eqn_wftmat}导出曲率表达式，
这再次说明式\eqref{chfb:eqn_wftmat}
（也就是式\eqref{chfb:eqn_wform-trans}）是微分几何中非常基本的重要公式。
对\eqref{chfb:eqn_wftmat}第一式求一次外微分，有
\begin{equation}\label{chfb:eqn_tmp-dwa1}
    \d (A \tilde{\omega})= \d \circ\d A + \d (\omega A)
    \ \Rightarrow \
    (\d A)\wedge \tilde{\omega}+ A \d \tilde{\omega}
    =(\d \omega) A - \omega\wedge \d A .
\end{equation}
其中矩阵之间的外积“$\wedge$”表示矩阵在相乘时，矩阵元素的积是外积。
我们以二维方阵为例来说明这个外积；设二维方阵$A$、$B$的矩阵元是1-形式场，有
\begin{align*}
	\begin{pmatrix}
		a_{11} & a_{12} \\ a_{21} & a_{22}
	\end{pmatrix} \wedge
	\begin{pmatrix}
		b_{11} & b_{12} \\ b_{21} & b_{22}
	\end{pmatrix}
	=\begin{pmatrix}
		a_{11}\wedge b_{11} + a_{12}\wedge b_{21} & a_{11}\wedge b_{12} + a_{12}\wedge b_{22} \\
		a_{21}\wedge b_{11} + a_{22}\wedge b_{21} & a_{21}\wedge b_{12} + a_{22}\wedge b_{22}
	\end{pmatrix}.
\end{align*}
将\eqref{chfb:eqn_wftmat}第一式改写为$\d A = A \tilde{\omega}-\omega A $，
然后代入\eqref{chfb:eqn_tmp-dwa1}，有
\begin{equation}\label{chfb:eqn_ATheta}
    A (\tilde{\omega}\wedge \tilde{\omega} + \d \tilde{\omega} )
    =(\d \omega +  \omega\wedge\omega )A .
\end{equation}

\begin{definition}
    $\Omega := \d \omega +  \omega\wedge\omega$ 叫作矢量丛联络$\nabla$在
    局部开集$U$上的\textbf{曲率形式}。
\end{definition}

$\Omega$是2-形式场。这样，式\eqref{chfb:eqn_ATheta}可记成：
\begin{equation}\label{chfb:eqn_TATA}
    \widetilde{\Omega} =A^{-1} \Omega A .
\end{equation}
这是曲率形式在局部标架变换时的变换公式。
上式是齐次的，即等号右端没有与$\Omega$无关的常数项；
而$\omega$的变换公式\eqref{chfb:eqn_wftmat}不是齐次的，
有非齐次项$A^{-1} \d A$。

设$X$、$Y$是$M$上任意两个切矢量场，则由曲率形式$\Omega$定义了从$\Gamma(E)$到$\Gamma(E)$的线性变换$R(X, Y)$。
$\forall p \in U\subset M$ 任取两个切矢量 $X, Y \in$ $T_p(M)$，则用曲率形式$\Omega$可定义从纤维$\pi^{-1}(p)$到
自身的线性变换$R(X, Y)$。它的定义如下：设 $\xi \in \pi^{-1}(p)$，它用矢量丛$E$在$U$上
的局部标架场$s_\alpha=\left(s_1, \cdots, s_q\right)$可表成
$\xi = \xi^\alpha s_{\alpha}$，其中$\xi^\alpha \in \mathbb{R}$。则命
\begin{equation}\label{chfb:eqn_CurvatureO}
    R(X, Y) \xi :=\sum\nolimits_{\alpha, \beta=1}^q \xi^\alpha
    \Omega_{\alpha}^\beta(X, Y) s_{\beta}|_p .
\end{equation}
每个$\Omega^\beta_{\alpha}$都是$U$上的2-形式；它作用于$X,Y\in T_pM$，
得到标量$\Omega^\beta_{\alpha}(X,Y)$。
这些标量组成$R(X,Y):E_p\to E_p$在标架$\{s_\alpha\}$下的矩阵，
并按式\eqref{chfb:eqn_TATA}作相似变换。
因此，式\eqref{chfb:eqn_CurvatureO}所定义的$R(X, Y)$是与局部标架选取无关的、从$\pi^{-1}(p)$到自身的线性变换。

如果$X$、$Y$是光滑流形$M$上的两个光滑切矢量场，则$R(X, Y)$是$\Gamma(E)$上的线性算子，
它的定义是：对任意的 $\xi \in \Gamma(E)$，及$p \in  M$，
\begin{equation}\label{chfb:eqn_RppR}
    \bigl(R(X, Y) \xi \bigr)(p):=R\left(X_p, Y_p\right) \xi_p .
\end{equation}

我们把$R(X, Y)$称为联络$\nabla$的\textbf{曲率算子}。

\begin{theorem}
    设$(E,M,\pi)$是矢量丛，$\forall X, Y\in \mathfrak{X}(M)$，$\forall \xi \in \Gamma(E)$，有
    \begin{equation}\label{chfb:eqn_Riemann-Curvature}
       R(X, Y)\xi=\nabla_X \nabla_Y\xi- \nabla_Y \nabla_X\xi-\nabla_{[X, Y]}\xi .
    \end{equation}
\end{theorem}

\begin{proof}
设$U$是$M$的局部开子集。因为外微分、联络具有局部属性，故曲率算子也是局部算子；
所以只要考虑式\eqref{chfb:eqn_Riemann-Curvature}两边分别在局部截面上的作用就行了。
设截面$\xi \in \Gamma(E|_U)$的局部表示为：$\xi =\sum_{\alpha=1}^q \xi^\alpha s_\alpha$，
其中$\xi^\alpha \in C^\infty(U)$。则有
\begin{align*}
    \nabla_X \xi= & \left(X (\xi^\alpha )+ X^i \xi^\beta \Gamma^\alpha_{i \beta} \right) s_\alpha
    = \bigl(X (\xi^\alpha )+ \xi^\beta \omega_{\beta}^\alpha(X) \bigr) s_\alpha . \\
    \nabla_Y \nabla_X \xi= & \Bigl(
    Y \bigl(X (\xi^\alpha )+ \xi^\beta \omega_{\beta}^\alpha(X) \bigr)+
    \bigl(X (\xi^\gamma )+ \xi^\beta \omega_{\beta}^\gamma(X) \bigr)
    \omega_{\gamma}^\alpha(Y)
    \Bigr) s_\alpha .
\end{align*}
利用 $1$-形式的外微分恒等式 $\d\omega(X,Y)=X(\omega(Y))-Y(\omega(X))-\omega([X,Y])$，得

\begin{align*}
    \nabla_X \nabla_Y \xi & -\nabla_Y \nabla_X \xi =
    \Bigl(
     X \bigl(Y (\xi^\alpha )+ \xi^\beta \omega_{\beta}^\alpha(Y) \bigr)+
     \bigl(Y (\xi^\gamma )+ \xi^\beta \omega_{\beta}^\gamma(Y) \bigr)
     \omega_{\gamma}^\alpha(X) \\
    &-Y \bigl(X (\xi^\alpha )+ \xi^\beta \omega_{\beta}^\alpha(X) \bigr)-
    \bigl(X (\xi^\gamma )+ \xi^\beta \omega_{\beta}^\gamma(X) \bigr)
    \omega_{\gamma}^\alpha(Y)     \Bigr)s_\alpha \\
   =&\Bigl( \left[X,Y\right](\xi^\alpha)
   +\xi^\beta \bigl( \d \omega_{\beta}^\alpha (X,Y)
   +\omega_{\beta}^\alpha([X,Y])
   + (\omega_{\gamma}^\alpha \wedge \omega_{\beta}^\gamma)(X,Y)
    \bigr)  \Bigr)s_\alpha  \\
   =&\nabla_{[X,Y]} \xi + \sum\nolimits_{\alpha,\beta=1}^{q}
   \xi^\beta \Omega_{\beta}^\alpha (X,Y) s_\alpha .
\end{align*}
由式\eqref{chfb:eqn_CurvatureO}可知定理得证。
\end{proof}

\begin{theorem}\label{chfb:thm_RP}
    有以下性质：
    (1) $R(X, Y)=-R(Y, X)$；
    (2) $R(f X, Y)=f \cdot R(X, Y)$；
    (3) $R(X, Y)(f \xi)=f \cdot\bigl(R(X, Y) \xi\bigr)$。
    其中$X, Y \in \Gamma(TM)$，$ f \in C^{\infty}(M)$，$ \xi \in \Gamma(E)$。
\end{theorem}
\begin{proof}
    证明过程与\S\ref{sec:appendix-curvature}类似，请读者补齐。
\end{proof}

\begin{theorem}\label{chfb:thm_2ndBianchi}
    曲率形式$\Omega$满足第二 Bianchi 恒等式：
    $ \d \Omega = \Omega \wedge \omega - \omega \wedge \Omega $。
\end{theorem}
\begin{proof}
对 $\Omega=\d \omega + \omega \wedge \omega$ 的两边求外微分得到
\begin{align*}
    \d \Omega & =\d \omega \wedge \omega- \omega \wedge \d \omega
     =(\Omega-\omega \wedge \omega) \wedge \omega-\omega \wedge(\Omega-\omega \wedge \omega) \\
    & =\Omega \wedge \omega-\omega \wedge \Omega .
\end{align*}
这里的楔积包含矩阵乘法。虽然标量 $2$-形式与 $1$-形式的楔积可交换，矩阵值形式 $\Omega$、$\omega$ 一般不可交换。
\end{proof}

\section{一般纤维丛}\label{chfb:sec_fb}
矢量丛的纤维是矢量空间，矢量丛是一种特殊的纤维丛。
一般纤维丛的纤维型是一个有Lie 变换群（见\S\ref{chlg:sec_Lie-transformation-group}）
左（或右）作用在其上的光滑流形。

\begin{definition}
    设$E$、$ M$、$ F$是光滑流形，$G$ 是左（或右）作用在流形$F$上的Lie 变换群。
    若有光滑满映射$\pi: E \rightarrow M$，并且下列条件成立：

    (1) $M$有一个开覆盖$\left\{U_\alpha: \alpha \in I\right\}$，
    并且对每个$\alpha \in I$有微分同胚
    \begin{equation*}
        \psi_\alpha: \pi^{-1}(U_\alpha) \to U_\alpha \times F,
        \quad \text{使得}\
        \pi \circ \psi_\alpha^{-1}(p, f)=p, \quad \forall p \in U_\alpha, f \in F
        \ \text{成立}.
    \end{equation*}

    (2) 对每一固定点 $p \in U_\alpha$，记
    $\psi_{\alpha, p}(f) := \psi_\alpha(p, f), \ \forall f \in F$，
    则 $\psi_{\alpha, p}: \pi^{-1}(p)\to F$是微分同胚；
    并且当$p \in U_\alpha \cap U_\beta$时，微分同胚$\psi_{\alpha, p}\circ \psi_{\beta, p}^{-1}:
    F \rightarrow F$是$G$的一个元素$g_{\alpha \beta}^p\in G$在$F$上的左（或右）作用，即
    $\psi_{\alpha, p} \circ \psi_{\beta, p}^{-1}(f)=g^p_{\alpha \beta} \cdot f,
    \ \forall f \in F$。

    (3) 当 $U_\alpha \cap U_\beta \neq \varnothing$ 时，
    映射 $g_{\alpha \beta}: U_\alpha \cap U_\beta \rightarrow G$ 是光滑的。

    则称 $(E, M, \pi, F, G)$ 为光滑流形 $M$ 上的一个\textbf{纤维丛}；
    $E$称为\textbf{丛空间}；$M$ 称为\textbf{底空间}；$F$ 称为\textbf{纤维型}；
    $\pi$称为\textbf{丛投影}；
    条件(1)中的映射称为\textbf{局部平凡化}。
\end{definition}

上述定义也可直接用局部平凡化表述：$\pi^{-1}(U_\alpha)\cong U_\alpha\times F$，
两组平凡化在交叠处由$(p,f)\mapsto(p,g_{\alpha\beta}(p)f)$联系。
底流形的坐标只用于检验这些映射的光滑性；丛本身由过渡函数及其相容条件确定。

一般纤维丛也有类似于定理\ref{chfb:thm_vector-existence}的存在性定理，叙述如下。

\begin{theorem}\label{chfb:thm_fb-existence}
    设 $M$ 是一个 $m$ 维光滑流形，$\left\{U_\alpha: \alpha \in I\right\}$是 $M$ 的一个开覆盖；
    $G$是左作用在光滑流形$F$上的Lie 变换群。
    如果对于任意一对指标 $\alpha, \beta \in I$，在 $U_\alpha \cap U_\beta \neq \varnothing$ 时，
    都指定了一个光滑映射 $g_{\alpha \beta}: U_\alpha \cap U_\beta \rightarrow G$，
    使得它们满足命题\ref{chfb:thm_trans-func}中的条件 (1)和 (2)
    （需将命题\ref{chfb:thm_trans-func}中的$GL(q)$换成Lie 变换群$G$）。
    那么存在一个纤维丛$(E,M,\pi,F,G)$，它以$F$为纤维型，并且它的过渡函数族
    恰好是 $\left\{g_{\alpha \beta}: U_\alpha \cap U_\beta \rightarrow G\right\}$ 。
\end{theorem}

\begin{proof}
    在不交并$\coprod_\alpha(U_\alpha\times F)$上规定
    $(p,f)_\beta\sim(p,g_{\alpha\beta}(p)f)_\alpha$。
    过渡函数的单位条件和余循环条件保证这是等价关系。
    以商空间为$E$，定义$\pi([(p,f)_\alpha])=p$。
    每个等价类在固定的$U_\alpha\times F$内有唯一代表，故得到
    $\psi_\alpha:\pi^{-1}(U_\alpha)\to U_\alpha\times F$。
    商映射是开映射，因而这些映射是同胚；其坐标变换正是
    $(p,f)\mapsto(p,g_{\alpha\beta}(p)f)$，故给出光滑结构。
    不同纤维上的点由底空间的开集分离，同一纤维上的点由一个局部平凡化分离，
    所以$E$是Hausdorff的；取底空间的可数子覆盖即可得到可数基。
    因此$E$满足纤维丛的全部条件。取$F=\mathbb R^q$及线性作用，便得到矢量丛的情形。
\end{proof}

\begin{definition}
    设$(E, M, \pi, F, G)$是一纤维丛，若$F=G$，
    且$G$在$F$上的左（或右）作用就是$G$在它自身上的左（或右）移动（左、右移动自然是自由且可迁的）；
    则称此丛为$M$上以$G$为结构群的\textbf{G-主丛}，记为 $(E, M, \pi, G)$。
\end{definition}

\begin{definition}\label{chfb:def_associated}
    设$(E, M, \pi, F, G)$和$(\widetilde{E}, M, \tilde{\pi}, \widetilde{F}, G)$是同一个光滑流形$M$上两个纤维丛。
    若存在$M$的一个开覆盖 $\left\{U_\alpha\right\}$，使得两个纤维丛有相同的过渡函数族
    \[
        \left\{g_{\alpha \beta}:U_\alpha \cap U_\beta \rightarrow G\right\},
    \]
    则称这两个纤维丛是\textbf{相配的}（associated）。
\end{definition}

纤维丛的相配关系是一个等价关系，具有自反性、对称性、传递性。

\begin{proposition}
    每一个纤维丛 $(E, M, \pi, F, G)$ 都有与它相配的$G$-主丛。
\end{proposition}
\begin{proof}
    设纤维丛$\pi: E \rightarrow M$的过渡函数族是$\left\{g_{\alpha \beta}: U_\alpha \cap U_\beta \rightarrow G\right\}$。
    $G$在它自身上的左移动可看作$G$左作用于自身的Lie 变换群，并且是自由的、可迁的；
    则根据定理\ref{chfb:thm_fb-existence}可知存在纤维丛
    $(\widetilde{E}, M, \tilde{\pi},G, G)$以$\left\{g_{\alpha \beta}: U_\alpha \cap U_\beta \rightarrow G\right\}$为
    过渡函数族，这是与原纤维丛$(E, M, \pi, F, G)$相配的$G$-主丛。
\end{proof}

这个命题再次强调了过渡函数族对于一般纤维丛的核心作用，
并且在结构群为$G$的纤维丛的相配类中确有一个$G$-主丛作为它的代表。

曲线提升一般不唯一，联络则通过指定其切向方向来选出提升。
设$\pi:E\to M$为光滑纤维丛，$V_u:=\ker\pi_{*u}$是沿纤维的\textbf{垂直子空间}。
光滑选取补空间$H_u$，使
\eq{T_uE=V_u\oplus H_u,}
称为给定一个\textbf{水平分布}，也称\textbf{Ehresmann 联络}。这里“水平”指所选的补空间，不涉及度规或正交性。
因$\pi_{*u}|_{H_u}:H_u\to T_{\pi(u)}M$是线性同构，每个$X\in T_{\pi(u)}M$恰有一个$X_u^H\in H_u$满足$\pi_{*u}X_u^H=X$，称为$X$在$u$处的\textbf{水平提升}。

\begin{definition}\label{chfb:def_horizontal-lift}
    设纤维丛$\pi:E\to M$已给定水平分布$H$。曲线$\gamma:I\to M$的提升$\widetilde\gamma:I\to E$若满足
    \eq{\dot{\widetilde\gamma}(t)\in H_{\widetilde\gamma(t)},}
    则称为$\gamma$的\textbf{水平提升}(horizontal lift)。等价地，
    \eq{\dot{\widetilde\gamma}(t)=\bigl(\dot\gamma(t)\bigr)^H_{\widetilde\gamma(t)}.}
\end{definition}

在局部丛坐标$(x^i,y^\alpha)$中，水平分布可唯一写成
\eq{H=\operatorname{span}\left\{\p_{x^i}+A_i^\alpha(x,y)\p_{y^\alpha}\right\}.}
故$\widetilde\gamma(t)=(x(t),y(t))$水平当且仅当
\eq{\dot y^\alpha=A_i^\alpha(x(t),y(t))\dot x^i.}
由常微分方程的存在唯一性，指定$\widetilde\gamma(t_0)\in E_{\gamma(t_0)}$后，水平提升局部存在且唯一。
一般水平分布不保证提升在整个参数区间存在。例如积丛$\R^2\to\R$、$(x,y)\mapsto x$上取$H=\operatorname{span}\{\p_x+y^2\p_y\}$，$\gamma(t)=t$从$(0,1)$出发的水平提升为$(t,(1-t)^{-1})$，不能延至$t=1$。

\begin{example}[水平提升与平行移动]
    矢量丛联络$\nabla$自然给出水平分布。在坐标$x^i$及局部丛标架$s_\alpha$中，记$\nabla_{\p_{x^i}}s_\alpha=\Gamma^\beta_{i\alpha}s_\beta$，令$v=y^\alpha s_\alpha(x)$，则
    \eq{\label{chfb:eq_horizontal-connection}
        H_{(x,y)}=\operatorname{span}\left\{
        \p_{x^i}-\Gamma^\beta_{i\alpha}(x)y^\alpha\p_{y^\beta}\right\}.}
    联络的换基律使这些子空间在交叠处一致。因而$v(t)$作为$E$中的曲线水平，当且仅当
    \eq{\dot v^\beta+\Gamma^\beta_{i\alpha}(\gamma(t))\dot x^i v^\alpha=0,}
    即式\eqref{chfb:eq_parallel-lift}中的$D v/\d t=0$。
    因此，\textbf{平行矢量场就是底曲线在矢量丛中的水平提升}；平行移动映射$P_\gamma$则记录该提升的终点。
    这里方程对纤维坐标线性，故给定初始矢量后，提升在底曲线的整个紧参数区间上存在。
\end{example}

对右$G$-主丛$P$，主丛联络要求水平分布还满足$(R_g)_*H_u=H_{ug}$。
这等价于给定$\mathfrak g$值一形式$\omega$，满足
$\omega(A^\#)=A$及$R_g^*\omega=\operatorname{Ad}_{g^{-1}}\omega$，
其中$A^\#_u=\left.\frac{\d}{\d t}\right|_0u\exp(tA)$，且$H_u=\ker\omega_u$。
此时曲线水平的条件就是$\omega(\dot{\widetilde\gamma})=0$。
相配矢量丛的截面可表示为等变函数$f:P\to V$；沿底空间矢量的水平提升对$f$求导，
便得到满足线性性和Leibniz法则的矢量丛联络。

\section{标架丛}\label{chfb:sec_frame-bundles}

标架丛不是矢量丛，它是$GL(q)$-主丛，它是矢量丛的相配丛。
相配丛是相互的，在纤维丛理论中，更多时候称矢量丛是$GL(q)$-标架丛的相配丛。

给定一个秩为$q$的矢量丛$(E,M,\pi)$，我们从矢量丛$E$出发构造光滑流形$M$上的标架丛；
它在局部上是$M$的开子集与一般线性群$GL(q)$的乘积空间，并且与矢量丛$E$共享同一个过渡函数族。
具体作法如下。

设矢量丛$(E,M,\pi)$的局部平凡化结构是$\left\{\left(U_{\alpha}, \psi_{\alpha}\right) ;\  \alpha \in I\right\}$，
其中$\left(U_{\alpha} ; x_{\alpha}^{i}\right)$是$M$的局部坐标系。
用$\left\{e_{i} ;\ 1 \leqslant i \leqslant q\right\}$表示实矢量空间$\mathbb{R}^{q}$的标准基底，并令
\begin{equation}
    s_{i}^{(\alpha)}(p)=\psi_{\alpha}^{-1}\left(p, e_{i}\right), \qquad p \in U_{\alpha} .
\end{equation}
则$s^{(\alpha)}=\bigl(s_{1}^{(\alpha)}, \cdots, s_{q}^{(\alpha)}\bigr)$是矢量丛$E$在$U_{\alpha}$上的局部标架场。
$\forall p \in U_{\alpha}$，用$F(p)$（“$F$”是Frame的首字母）表示由矢量空间$\pi^{-1}(p)$中所有基底构成的集合，
则在$GL(q)$和$F(p)$之间有双射关系。验证如下，令$A \in GL(q)$所对应的标架是
\begin{equation}\label{chfb:eqn_F-comp}
    f(p)  =\bigl(f_{1}(p), \cdots, f_{q}(p)\bigr)
    =\left(s_{1}^{(\alpha)}(p), \cdots, s_{q}^{(\alpha)}(p)\right) \cdot A = s^{(\alpha)}(p) \cdot A,
\end{equation}
注意$A \in GL(q)$对标架是右作用，对矢量分量是左作用。上式即为
\begin{equation}
    f_{i}(p)=A_{i}^{j} \psi_{\alpha}^{-1}(p, e_{j})
    =\psi_{\alpha, p}^{-1} (A_{i}^{j} e_{j}), \qquad 1 \leqslant i \leqslant q .
\end{equation}
因$\psi_{\alpha, p}^{-1}$是线性同构，故矩阵$A$可移进移出。这说明了$GL(q)$和$F(p)$是双射。
\begin{equation}
    \text{令}\quad  F(E) := \bigcup_{p \in M} F(p) ,
\end{equation}
并且定义映射$\tilde{\pi}: F(E) \rightarrow M$使得
\begin{equation}
    \tilde{\pi}\bigl(F(p)\bigr)=\{p\} , \quad \forall p \in M.
\end{equation}
下面借助矢量丛$E$的局部平凡化，引入空间$F(E)$的局部平凡化结构。
事实上，$\forall \alpha \in I$，定义映射
$\varphi_{\alpha}: \tilde{\pi}^{-1}(U_{\alpha})\to U_{\alpha} \times GL(q) $，使得
\begin{equation}
    \varphi_{\alpha}^{-1}(p, A) := s^{(\alpha)}(p)
     \cdot A, \qquad \forall(p, A) \in U_{\alpha} \times GL(q) .
\end{equation}
容易看出，这样定义的映射$\varphi_{\alpha}$是从$\tilde{\pi}^{-1}\left(U_{\alpha}\right)$到$U_{\alpha} \times GL(q)$上的双射。
借助映射$\varphi_{\alpha}$可以在$F(E)$上引进拓扑结构和光滑结构，使得每一个$\varphi_{\alpha}$是微分同胚，
并且$\tilde{\pi}: F(E) \rightarrow M$是光滑映射；具体作法可以参照例\ref{chfb:exm_ctb}。
特别地，通过映射$\varphi_{\alpha}: \tilde{\pi}^{-1}(U_{\alpha}) \rightarrow U_{\alpha} \times GL(q)$以及$M$在
坐标邻域$U_{\alpha}$上的坐标映射可以给出$F(E)$在$\tilde{\pi}^{-1}(U_{\alpha})$上的局部坐标系，
记为$\left(\tilde{\pi}^{-1}\left(U_{\alpha}\right) ; x_{\alpha}^{i}, A_{k}^{j} \right)$。

从上面的构造可以看出，$F(E)$在局部上是乘积空间$U_{\alpha} \times GL(q)$，
它在每一点$p \in M$处的纤维是$\tilde{\pi}^{-1}(p)=F(p)$，后者与$GL(q)$是等同的。
换言之，$F(E)$是将一族乘积空间$U_{\alpha} \times GL(q), \alpha \in I$，沿底流形$M$上
同一点$p \in U_{\alpha} \cap U_{\beta}$的纤维$\{p\} \times GL(q) \subset U_{\alpha} \times GL(q)$
和$\{p\} \times GL(q) \subset U_{\beta} \times GL(q)$按照一定规则粘合起来的结果。
具体地说，设$p \in U_{\alpha} \cap U_{\beta}$，$ A, B \in GL(q)$，那么
\begin{equation}\label{chfb:eqn_phiequ}
    \varphi_{\alpha}^{-1}(p, A)=\varphi_{\beta}^{-1}(p, B)
    \quad \Leftrightarrow \quad
    A=\varphi_{\alpha, p} \circ \varphi_{\beta, p}^{-1}(B) ,
\end{equation}
当且仅当 （其中$1 \leqslant i \leqslant q $）
\begin{equation}
    s^{(\alpha)}(p) \cdot A=s^{(\beta)}(p) \cdot B
    \quad \Leftrightarrow \quad
    \psi_{\alpha}^{-1}(p, A_{i}^{j} e_{j})
    =\psi_{\beta}^{-1}(p, B_{i}^{j} e_{j}).
\end{equation}
上式等价于 （注意$\psi_{\alpha, p} \circ \psi_{\beta, p}^{-1}$是线性同构，矩阵$B$可移进移出）
\begin{equation}
    A_{i}^{j} e_{j}=\psi_{\alpha, p} \circ \psi_{\beta, p}^{-1}(B_{i}^{j}e_{j})
    =B_{i}^{k}\ \psi_{\alpha, p} \circ \psi_{\beta, p}^{-1}(e_{k})
    =\bigl(g_{\alpha \beta}^p\bigr)_{k}^{j} B_{i}^{k} e_{j} .
\end{equation}
即有
\begin{equation}\label{chfb:eqn_agb}
    A=g_{\alpha \beta}^p \cdot B \quad \Leftrightarrow \quad
    A_{i}^{j}  = \bigl(g_{\alpha \beta}^p\bigr)_{k}^{j} B_{i}^{k}.
\end{equation}
将上式与式\eqref{chfb:eqn_phiequ}相比较得知：
\begin{equation}
    \varphi_{\alpha, p} \circ \varphi_{\beta, p}^{-1}=g_{\alpha \beta}^p
    =\psi_{\alpha, p} \circ \psi_{\beta, p}^{-1}.
\end{equation}
通过上面论述可知：标架丛$\tilde{\pi}: F(E) \rightarrow M$和矢量丛$\pi: E \rightarrow M$共享同一个过渡函数族
$\left\{g_{\alpha \beta}: U_{\alpha} \cap U_{\beta} \rightarrow GL(q)\right\}$；它们的底流形也相同。
矢量丛的纤维型是$\mathbb{R}^q$；标架丛的纤维型是$GL(q)$，故它是$GL(q)$-主丛。
依定义\ref{chfb:def_associated}称纤维丛$\tilde{\pi}: F(E) \rightarrow M$是
与矢量丛$\pi: E \rightarrow$ $M$相配的\textbf{标架丛}，
记作$\bigl(F(E), M, \tilde{\pi}\bigr)$，或$\tilde{\pi}: F(E) \rightarrow M$，或$F(E)$。

对前面的叙述做一个小结：
设$\left\{\left(U_{\alpha}, \psi_{\alpha}\right) ; \alpha \in I\right\}$是矢量丛$\pi: E \rightarrow M$的
局部平凡化结构，使得$\left(U_{\alpha} ; x_{\alpha}^{i}\right)$是底流形$M$的局部坐标系，
则$\left\{\left(U_{\alpha}, \varphi_{\alpha}\right) ; \alpha \in I\right\}$是与$E$相配的
标架丛$\tilde{\pi}: F(E) \rightarrow M$的一个局部平凡化结构，
且通过局部平凡化映射$\varphi_{\alpha}: \tilde{\pi}\left(U_{\alpha}\right) \rightarrow U_{\alpha} \times GL(q)$给出了
丛空间$F(E)$上的局部坐标系$\left(\tilde{\pi}^{-1}\left(U_{\alpha}\right) ; x_{\alpha}^{i}, A_{a}^{b}\right)$。

下面由矢量丛联络 $\nabla$ 构造标架丛上的联络形式。活动标架 $f$ 视为拉回丛 $\tilde\pi^*E$ 的标架，其诱导协变导数记为 $\nabla^F$。
任取$\alpha \in I$，然后令其固定不变；有
$ \nabla s_{a}^{(\alpha)}=\omega_{a}^{(\alpha) b}\otimes s_{b}^{(\alpha)}$、
$ \omega_{a}^{(\alpha) b}=\Gamma^{(\alpha) b}_{i a} \d x_{\alpha}^{i} $
成立，其中$\omega_{a}^{(\alpha) b}$是联络$\nabla$在局部标架场$\{s_{a}^{(\alpha)}\}$下的联络形式。
采用式\eqref{chfb:eqn_F-comp}的记号，则活动标架$f=s^{(\alpha)} \cdot A$的联络定义为：
\begin{equation}
\begin{aligned}
    \nabla^F f & :=\left(\nabla s^{(\alpha)}\right) \cdot A + s^{(\alpha)} \cdot \d A
    =s^{(\alpha)} \cdot \left(\omega^{(\alpha)} \cdot A+\d A\right)  \\
    & =f \cdot A^{-1}\left(\d A+\omega^{(\alpha)} \cdot A\right)
    =f \cdot \theta .
\end{aligned}\end{equation}
这里矩阵记法中的$f\cdot\theta$表示各列$\sum_b\theta^b_a\otimes f_b$，
$s^{(\alpha)}\cdot\d A$也按同样方式理解；底流形上的形式均拉回至$F(E)$。
其中$\theta$表达式见\eqref{chfb:eqn_lt}。之后，将$\nabla^F$简记为$\nabla$。
\begin{equation}\label{chfb:eqn_lt}
    \theta_{a}^{b}\! =\! \left(A^{-1}\right)_{c}^{b}
    \left(\d A_{a}^{c}+\omega_{d}^{(\alpha) c} A_{a}^{d} \right)
    =\left(A^{-1}\right)_{c}^{b}\left(\d A_{a}^{c}+A_{a}^{d}
     \Gamma^{(\alpha) c}_{i d} \d x_{\alpha}^{i}\right) .
\end{equation}

\begin{theorem}
由式\eqref{chfb:eqn_lt}定义的$\theta_{a}^{b} (1 \leqslant a,b \leqslant q)$是大范围
定义在丛空间$F(E)$上的$q^{2}$个1-形式，即$\theta=(\theta_{a}^{b})$是大范围定义在$F(E)$上
的$\mathfrak{g l}(q)$-值的1-形式，其中$\mathfrak{gl}(q)$是一般线性群$GL(q)$的Lie 代数。
\end{theorem}

\begin{proof}
设$\left(U_{\beta} ; x_{\beta}^{i}\right)$是$M$的另一局部坐标系，
$\left(\tilde{\pi}^{-1}\left(U_{\beta}\right) ; x_{\beta}^{i}, B_{a}^{b}\right)$是
在$F(E)$上对应的局部坐标系。若$U_{\alpha} \cap U_{\beta} \neq \varnothing$，
则有坐标变换$A=g_{\alpha \beta}^p \cdot B$（参看式\eqref{chfb:eqn_agb}），
其中$g_{\alpha \beta}: U_{\alpha} \cap U_{\beta} \rightarrow GL(q)$是过渡函数。
对\eqref{chfb:eqn_agb}取外微分，有
\begin{equation}\label{chfb:eqn_tmp-dabc}
    \d A_{a}^{b}=\d \left(g_{\alpha \beta}\right)_{c}^{b} \cdot
    B_{a}^{c}+\left(g_{\alpha \beta}\right)_{c}^{b} \cdot \d B_{a}^{c}.
\end{equation}

另一方面，联络形式遵循如下的变换公式（即，式\eqref{chfb:eqn_wform-trans}）
\begin{equation}\label{chfb:eqn_tmp-gb}
    \left(g_{\alpha \beta}\right)_{c}^{a} \omega_{b}^{(\beta) c}=\d
    \left(g_{\alpha\beta}\right)_{b}^{a}+
    \omega_{c}^{(\alpha) a} \left(g_{\alpha \beta}\right)_{b}^{c}.
\end{equation}
将式\eqref{chfb:eqn_tmp-gb}代入式\eqref{chfb:eqn_tmp-dabc}得到
\begin{align*}
    &\d A_{a}^{b}+A_{a}^{c}  \omega_{c}^{(\alpha) b}
    =\left(g_{\alpha \beta}\right)_{c}^{b}\left(\d B_{a}^{c}
    +B_{a}^{d}  \omega_{d}^{(\beta) c}\right).
\quad \xRightarrow{\ref{chfb:eqn_agb}} \\
&\left(A^{-1}\right)_{d}^{b} \left(\d A_{a}^{d}
+A_{a}^{c}  \omega_{c}^{(\alpha) d}\right)
=\left(B^{-1}\right)_{d}^{b}\left(\d B_{a}^{d}
+B_{a}^{c}  \omega_{c}^{(\beta) d}\right) .
\end{align*}
上式表明定义与局部标架选取无关，故在大范围上有定义。
\end{proof}

丛 $E$ 未必有全局标架。拉回丛 $\tilde\pi^*E$ 则有典范标架：
在 $f=(f_1,\ldots,f_q)\in F(E)$ 处取 $f_1,\ldots,f_q$。
这里的标架属于拉回丛，不是 $F(E)$ 的切标架。

\paragraph{对偶标架的联络}
对偶标架$\{f^{i}\}$由$f^{i}(f_j)=\delta^i_j$确定。
对任意$X\in\mathfrak{X}(F(E))$，由$\nabla_X f_j=\theta^b_j(X)f_b$及对偶联络的定义，有
\[
    (\nabla_X f^{i})(f_j)
    =X\bigl(f^{i}(f_j)\bigr)-f^{i}(\nabla_X f_j)
    =-\theta^i_j(X).
\]
因此
\begin{equation}
    \nabla f^{i}=-\theta^i_j\otimes f^{j}.
\end{equation}
这与第\ref{appx:curvature}章切丛及其对偶联络的形式一致。

\subsection{切标架丛}

取 $E=TM$，即得\textbf{切标架丛} $F(TM)$。
切丛与切标架丛共享过渡函数族，但纤维型分别为 $\R^m$ 和 $GL(m)$。
第\ref{appx:curvature}章的联络形式、Cartan 结构方程和 Bianchi 恒等式均可在此框架下表述。

若 $M$ 具有正定度规，正交标架构成结构群为 $O(m)$ 的\textbf{正交切标架丛}，记为 $\mathcal O(TM)$。
选定定向后，正向正交标架的结构群为 $SO(m)$。
四维 Lorentz 流形的正交标架丛结构群为 $O(1,3)$；同时选定时空定向和时间定向后，可约化为 $SO^+(1,3)$。

\section{矢量丛与规范场概述}\label{chfb:sec_G-connection}

规范势是主丛联络在局部截面下的拉回$A$，曲率为
$F=\d A+\tfrac12[A\wedge A]$；在矩阵表示中即$F=\d A+A\wedge A$。
更换局部截面时，$A'=g^{-1}Ag+g^{-1}\d g$，而$F'=g^{-1}Fg$。
以下用相配矢量丛的局部标架来描述这些量。

设底流形$M$是$m$维的光滑实流形，$M$上植有$G$-标架丛（$G$是$GL(q,\mathbb{C})$的子群）。
设群$G$的Lie 代数是$\mathfrak{g}$（Lie 代数的维数是$r$），
$\mathfrak{g}$的基为$\{T_a\mid a=1,\cdots,r\}$（$T_a$是$q$维方阵）。
若 $G$ 半单，其 Lie 代数的 Killing 型 $B$ 非退化；若 $G$ 紧致半单，则 $-B$ 正定。
它给出 $\mathfrak g$ 上的 $\operatorname{Ad}(G)$-不变双线性型，并诱导伴随丛的\textbf{丛度量}。
一般矢量丛 $E$ 的丛度量需由其纤维表示另行指定。

在开子集$U\subset M$、$V\subset M$上分别引入局部坐标$\{x^i\}$、$\{\tilde{x}^i\}$，并假定$U\cap V\neq \varnothing$；
在纤维型$\mathbb{C}^q$上引入局部标架$s_{\alpha}$（$1\leqslant \alpha \leqslant q$）；
故矢量丛有局部标架$\{\frac{\partial }{\partial x^i}; s_\alpha\}$。
在上述前提下，与矢量丛相配的$G$-标架丛（$G$-主丛）的联络形式$\theta$为：
\begin{equation}
	\theta_{a}^{b} = \left(a^{-1}\right)_{\beta}^{b}
	\left(\d a_{a}^{\beta}
	+\omega_{\alpha}^{\beta} a_{a}^{\alpha} \right) ,
	\qquad \text{其中}\  a_{a}^b\in {G}     \tag{\ref{chfb:eqn_lt}}
\end{equation}
其中$\omega_{\alpha}^\beta$是：
\begin{equation}\label{chfb:eqn_Gwform}
	\omega_{\alpha}^\beta=\sum\nolimits_{a=1}^{r} \sum\nolimits_{i=1}^{m}
	A^a_i(x) (T_a)_{\alpha}^\beta \d x^i;
	\ 1\leqslant \alpha , \beta \leqslant q,\ A^a_i(x)\in C^\infty(U).
\end{equation}
上式中$T_a$是结构群$G$的Lie 代数$\mathfrak{g}$的基矢矩阵，它们是$q$维常数方阵。

{我们省略了大量内容（$G$-主丛联络理论）才得到式\eqref{chfb:eqn_Gwform}。
	$G$-主丛理论中有定理表明在主丛联络中可以自然地引入结构群$G$的Lie 代数$\mathfrak{g}$，
	见式\eqref{chfb:eqn_Gwform}。若不使用$G$-主丛，只考虑矢量丛，
	则在矢量丛联络\ref{chfb:def_vb-conncection}中无法自然引入结构群$G$的Lie 代数$\mathfrak{g}$。
	我们借用$G$-主丛的目的就是要在式\eqref{chfb:eqn_Gwform}中自然引入$G$的Lie 代数$\mathfrak{g}$。}

当$G$-标架丛是$GL(m)$-切标架丛时，其结构群便是$GL(m)$；与切标架丛相配的矢量丛就是切丛。
此时$G$-标架丛联络\eqref{chfb:eqn_Gwform}就是定义\ref{def:connection-one-form}中的联络 $1$-形式（请读者验证之）。

由\S\ref{chfb:sec_vc}、\S\ref{chfb:sec_frame-bundles}内容可知由矢量丛联络可诱导出标架丛联络。
反之，若先给定$G$-标架丛联络（见上面两式），$G$-主丛联络理论表明：标架丛联络也可自然诱导出矢量丛联络；
两个丛的联络是相互唯一确定的，即式\eqref{chfb:eqn_Dsos}和式\eqref{chfb:eqn_lt}。
由$G$-标架丛联络\eqref{chfb:eqn_Gwform}诱导出的与之相配矢量丛的联络$\omega$自然
满足式\eqref{chfb:eqn_wform-trans}：
\begin{equation}
	\tilde{\omega}_{\alpha}^\sigma=(a^{-1})_{\beta}^{\sigma}
	\d a_{\alpha}^\beta+(a^{-1})_{\beta}^{\sigma}
	\omega_{\gamma}^\beta a_{\alpha}^\gamma ,
	\qquad \text{其中}\  a_{\gamma}^\beta\in {G}
	\tag{\ref{chfb:eqn_wform-trans}}
\end{equation}
至此，我们由$G$-标架丛联络得到了与之相配矢量丛的联络方阵$\omega$（见\eqref{chfb:eqn_Gwform}）。

\paragraph{矢量丛$G$型联络}
由式\eqref{chfb:eqn_Gwform}可得矢量丛联络系数表达式（见式\eqref{chfb:eqn_con-form}）：
\begin{equation}\label{chfb:eqn_GGamma}
	\Gamma^\beta_{i \alpha}=\omega_{\alpha}^\beta(\partial_i)
	=\sum\nolimits_{a=1}^{r} A^a_i(x) (T_a)_{\alpha}^\beta ;
	\quad 1\leqslant \alpha , \beta \leqslant q.
\end{equation}
把式\eqref{chfb:eqn_wform-trans}写成联络系数的形式
（$a_{\gamma}^\beta=[g_{VU}]_{\gamma}^\beta \in {G}$）：
\begin{equation}\label{chfb:eqn_GDG}
	\bigl[g_{VU}(x)  \bigr]_{\gamma}^\beta \tilde{\Gamma}_{j \alpha}^\gamma (\tilde{x})
	=\left(\frac{\partial \bigl[g_{VU}(x)  \bigr]_{\alpha}^\beta}{\partial x^i}
	+\bigl[g_{VU}(x)  \bigr]_{\alpha}^\gamma \Gamma^\beta_{i \gamma}(x)
	\right) \times \frac{\partial x^i }{\partial \tilde{x}^j}.
\end{equation}
满足变换关系\eqref{chfb:eqn_GDG}的式\eqref{chfb:eqn_GGamma}称为矢量丛的\textbf{$G$型联络}。

\begin{example}
	联络系数\eqref{chfb:eqn_GDG}的变换关系可以推广到一般情形。
\end{example}
联络系数\eqref{chfb:eqn_GDG}是在自然标架$\{\frac{\partial }{\partial x^i}\}$下的关系式。
我们将自然标架换成一般的切标架场$\{e_i\}$，
那么在变换$\tilde{s}_\alpha=a_{\alpha}^\gamma s_\gamma$、
$\tilde{e}_j=\mathcal{C}_{j}^i e_i$下，有
\begin{align*}
	\tilde{s}_\gamma \tilde{\Gamma}^\gamma_{j \alpha}
	=&  \nabla_{\tilde{e}_j} \tilde{s}_\alpha
	= \mathcal{C}_{j}^i \nabla_{e_i} \left(a_{\alpha}^\gamma s_\gamma\right)
	= \mathcal{C}_{j}^i \left( \left(\nabla_{e_i} a_{\alpha}^\gamma\right) s_\gamma
	+ a_{\alpha}^\gamma  \nabla_{e_i} s_\gamma\right) \\
	=& \mathcal{C}_{j}^i s_\beta \left({e_i}\left( a_{\alpha}^\beta\right)
	+ a_{\alpha}^\gamma \Gamma^\beta_{i \gamma} \right) .
\end{align*}
由上式可以得到
\begin{equation}\label{chfb:eqn_gfGDG}
	a_{\gamma}^\beta \tilde{\Gamma}^\gamma_{j \alpha} = \mathcal{C}_{j}^i
	\left({e_i}\left( a_{\alpha}^\beta\right)
	+ a_{\alpha}^\gamma \Gamma^\beta_{i \gamma} \right) .
\end{equation}
当把$\{e_i\}$取成自然标架$\{\frac{\partial }{\partial x^i}\}$时，
式\eqref{chfb:eqn_gfGDG}变成了式\eqref{chfb:eqn_GDG}。
\qed

\subsection{曲率相关内容的分量表述}

取$U$的自然标架$\{\frac{\partial }{\partial x^i}\}$、丛空间的标架$\{s_\alpha\}$；
对于$X=X^i \partial_i$、$Y=Y^j \partial_j$，可见$R(X, Y) s_\alpha \in \Gamma(E)$；
由\eqref{chfb:eqn_CurvatureO}、\eqref{chfb:eqn_RppR}得（其中$\xi = \xi^\alpha(x) s_{\alpha}$）
\begin{align*}
	s^{\gamma}\bigl(R(X,Y)\xi\bigr)
	=\xi^\alpha \Omega_{\alpha}^\beta(X, Y) s^{\gamma}(s_\beta)
	=\xi^\alpha \Omega_{\alpha}^\gamma (X, Y) .
\end{align*}
故有
\begin{subequations}
	\begin{align}
	R_{\alpha i j}^\beta := &
	\Omega_{\alpha}^\beta (\partial_i, \partial_j) =	s^{\beta}\bigl(R(\partial_i,\partial_j)s_\alpha\bigr) .   \\
	R =& R_{\alpha i j}^\beta \ s_\beta \otimes
	s^{\alpha}\otimes \d x^i \otimes \d x^j . 
	\end{align}
\end{subequations}
此处$R\in\Gamma(E\otimes E^*\otimes T^*M\otimes T^*M)$。
对$\lambda\in E_p^*$、$\xi\in E_p$及$X,Y\in T_pM$，其多重线性作用为
$R(\lambda;\xi,X,Y)=\lambda\bigl(R(X,Y)\xi\bigr)$。

\paragraph{坐标变换} 下面导出在变换$\tilde{s}_\alpha=a_{\alpha}^\gamma s_\gamma$、
$x^i\to \tilde{x}^i$下，曲率分量表达式是如何变换的。
我们从式\eqref{chfb:eqn_TATA}、\eqref{chfb:eqn_CurvatureO}、\eqref{chfb:eqn_RppR}出发进行探讨。
\begin{align*}
	\widetilde{\Omega} =A^{-1} \Omega A \quad \Rightarrow \quad
	\widetilde{\Omega}_{\alpha}^\beta
	\left( \frac{\partial }{\partial \tilde{x}^i}, \frac{\partial }{\partial \tilde{x}^j} \right)
	=\left(a^{-1}\right)_{\rho}^\beta a_{\alpha}^\sigma
	\Omega_{\sigma}^\rho
	\left( \frac{\partial x^k }{\partial \tilde{x}^i} \frac{\partial }{\partial x^k},
	\frac{\partial x^l }{\partial \tilde{x}^j} \frac{\partial }{\partial x^l} \right) .
\end{align*}
由此可以得到局部坐标变换下，曲率分量的变换表达式（$a\in G$）：
\begin{equation}
	\widetilde{R}^\beta_{\alpha ij}
	= R^\rho_{\sigma kl}
	\left(a^{-1}\right)_{\rho}^\beta a_{\alpha}^\sigma
	\frac{\partial x^k }{\partial \tilde{x}^i} \frac{\partial x^l }{\partial \tilde{x}^j} .
\end{equation}

\paragraph{曲率分量表达式} 由式\eqref{chfb:eqn_Riemann-Curvature}可以得到
（请与式\eqref{eq:curvature-components}比对）：
\begin{equation}\label{chfb:eqn_R-vector-bundle}
	R_{\alpha i j}^\beta
	= \frac{\partial}{\partial x^i} \Gamma^{\beta}_{j \alpha}
	-\frac{\partial}{\partial x^j} \Gamma^{\beta}_{i \alpha}
	+\Gamma^{\sigma}_{j \alpha} \Gamma^{\beta}_{i \sigma}
	-  \Gamma^{\sigma}_{i \alpha}\Gamma^{\beta}_{j \sigma} .
\end{equation}
上式中各分量指标的取值范围是：$ 1\leqslant i,j \leqslant m $，$ 1\leqslant \alpha,\beta,\sigma \leqslant q$。

\paragraph{曲率协变导数} 先计算局部自然标架上的导数：

\begin{align*}
	\nabla_{\partial_k} R = & \nabla_{\partial_k}\left(	R_{\alpha i j}^\beta
	s_\beta \otimes s^{\alpha}\otimes \d x^i \otimes \d x^j \right) \\
	=& \frac{\partial R_{\alpha i j}^\beta }{\partial x^k}
	s_\beta \otimes s^{\alpha}\otimes \d x^i \otimes \d x^j
	+ R_{\alpha i j}^\beta \left( \nabla_{\partial_k} s_\beta \right)
	\otimes s^{\alpha}\otimes \d x^i \otimes \d x^j \\
	& + R_{\alpha i j}^\beta s_\beta \otimes
	\left( \nabla_{\partial_k} s^{\alpha} \right) \otimes \d x^i \otimes \d x^j
	+ R_{\alpha i j}^\beta s_\beta \otimes s^{\alpha}\otimes
	\left( \nabla_{\partial_k} \d x^i \right) \otimes  \d x^j \\
	&+ R_{\alpha i j}^\beta s_\beta \otimes s^{\alpha}\otimes \d x^i
	\otimes \left( \nabla_{\partial_k} \d x^j \right) \\
	=&\left( \frac{\partial R_{\alpha i j}^\beta }{\partial x^k}
	+ R_{\alpha i j}^\rho \Gamma^\beta_{k \rho}
	- R_{\rho i j}^\beta \Gamma^\rho_{k \alpha}
	- R_{\alpha l j}^\beta \gamma^l_{k i}
	- R_{\alpha i l}^\beta \gamma^l_{k j}	\right)
	s_\beta \otimes s^{\alpha}\otimes \d x^i \otimes \d x^j .
\end{align*}
为对曲率张量求协变导数，另在 $TM$ 上选取仿射联络。这里 $\Gamma^\beta_{k \alpha}$ 是矢量丛 $E$ 的联络系数，$\gamma^l_{k i}$由$\nabla_{\partial_k}\partial_i=\gamma^l_{k i}\partial_l$定义，是切丛$TM$的仿射联络系数。我们记
\begin{equation}\label{chfb:eqn_RijDk}
	 R_{\alpha i j ; k}^\beta := \frac{\partial R_{\alpha i j}^\beta }{\partial x^k}
	+ R_{\alpha i j}^\rho \Gamma^\beta_{k \rho}
	- R_{\rho i j}^\beta \Gamma^\rho_{k \alpha}
	- R_{\alpha l j}^\beta \gamma^l_{k i}
	- R_{\alpha i l}^\beta \gamma^l_{k j} .
\end{equation}
对于任意的$ Z = Z^k \partial_k \in \mathfrak{X}(U)$，则有
\begin{equation}
	\nabla_{Z} R = Z^k R_{\alpha i j ; k}^\beta \
	s_\beta \otimes s^{\alpha}\otimes \d x^i \otimes \d x^j .
\end{equation}

\paragraph{Bianchi 微分恒等式分量表达式}
需要注意曲率方阵$\Omega = (\Omega_{\alpha}^\beta)$中的每个
矩阵元$\Omega_{\alpha}^\beta$都是$U\subset M$上的$2$形式场，
即$\Omega_{\alpha}^\beta= \frac{1}{2}
R_{\alpha ij}^\beta \d x^i \wedge \d x^j$，
其中$R_{\alpha i j}^\beta \in C^\infty(U)$。
将其代入定理\ref{chfb:thm_2ndBianchi}中
的式子$\d \Omega = \Omega \wedge \omega - \omega \wedge \Omega $，有
\begin{align*}
    0 = & 2(\d \Omega - \Omega \wedge \omega + \omega \wedge \Omega )\\
    =& \partial_k R^\beta_{\alpha ij} \d x^k \wedge \d x^i \wedge \d x^j
    - R^\beta_{\rho ij} \d x^i \wedge \d x^j \wedge \omega^\rho_\alpha
    + \omega^\beta_\rho \wedge R^\rho_{\alpha ij} \d x^i \wedge \d x^j \\
    \xlongequal{\ref{chfb:eqn_con-form}}&
    \left( R^\beta_{\alpha ij,k}
    - R^\beta_{\rho ij} \Gamma^\rho_{k \alpha}
    + \Gamma^\beta_{k \rho} R^\rho_{\alpha ij}
    \right)\d x^k \wedge \d x^i \wedge \d x^j \\
    \xlongequal{\ref{chfb:eqn_RijDk}} &
    \left(R^\beta_{\alpha i j ; k}
    + R^\beta_{\alpha l j} \gamma^l_{k i}
    + R^\beta_{\alpha i l} \gamma^l_{k j}
    \right)\d x^k \wedge \d x^i \wedge \d x^j \\
    \xlongequal[\ref{chfb:thm_RP}]{\text{定理}} &
    \left(R^\beta_{\alpha i j ; k}
    + R^\beta_{\alpha l j} \gamma^l_{k i}
    - R^\beta_{\alpha l i} \gamma^l_{k j}
    \right)\d x^k \wedge \d x^i \wedge \d x^j .
\end{align*}
利用外微分场的全反对称性，有

\begin{align*}
	0 = R^\beta_{\alpha i j ; k}
	+ {R^\beta_{\alpha l j} \gamma^l_{k i} }
	- {R^\beta_{\alpha l i} \gamma^l_{k j} }
	+ R^\beta_{\alpha  j k ; i}
	+ {R^\beta_{\alpha l i} \gamma^l_{j k} }
	- {R^\beta_{\alpha l j} \gamma^l_{i k} }
	+ R^\beta_{\alpha k i ; j}
	+ {R^\beta_{\alpha l k} \gamma^l_{i j} }
	- {R^\beta_{\alpha l k} \gamma^l_{j i} } .
\end{align*}
若仿射联络满足无挠条件（$\gamma^l_{j k}=\gamma^l_{k j}$），
则可得Bianchi 微分恒等式：
\begin{equation}
	R^\beta_{\alpha i j ; k} + R^\beta_{\alpha  j k ; i}
	 + R^\beta_{\alpha k i ; j} = 0
	 \quad \Leftrightarrow \quad
	 R^\beta_{\alpha [ij ; k]} = 0 .
\end{equation}

\subsection{平直时空规范场}

在经典场论中，需将$M$选为平直的四维闵氏流形$(\mathbb{R}^4_1,\eta)$（$\gamma$全是零）；
且不作坐标变换$x^i\to \tilde{x}^j$，只作变换$s_\alpha\to \tilde{s}_\beta$。
以下两例就在这样的假设下来阐述。

\begin{example}
	$U(1)$规范场。
\end{example}

取 $G=U(1)=\{\e^{\i\theta}:\theta\in\R\}$，其 Lie 代数为 $\mathfrak u(1)=\i\R$，选基 $T_1=\i$。

即令式\eqref{chfb:eqn_GDG}中的$g_{VU}(x)=\e^{\i \theta}$、$T_a=(\i )$，
则式\eqref{chfb:eqn_GDG}化为：
\begin{equation*}
	\e^{\i \theta} \tilde{\Gamma}_{j 1}^{1}
	=\frac{\partial \e^{\i \theta}}{\partial x^j}
	+\Gamma_{j 1}^1 \e^{\i \theta}
	\quad \xRightarrow[\ref{chfb:eqn_GGamma}]{\text{代入式}}\quad
	\tilde{A}_{j}(x) =\frac{\partial \theta}{\partial x^j}  +A_j(x)  .
\end{equation*}
这正是电磁规范势的规范变换公式；此时式\eqref{chfb:eqn_GGamma}中的$A^a_i$对应{电磁规范势}。

把式\eqref{chfb:eqn_GGamma}（$U(1)$群是可对易群）代入曲率公式\eqref{chfb:eqn_R-vector-bundle}，有
\begin{equation}\label{chfb:eqn_EM-strength}
	R_{1 i j}^1 = \i ({\partial_i} A_{j} -\partial_j A_i) .
\end{equation}
上式正是电磁场强公式$\i ^{-1} R_{1 i j}^1= F_{ij} = {\partial_i} A_{j} -\partial_j A_i$。
\qed

\begin{example}
	非对易规范场（结构群$G$是半单、紧致、非对易方阵群）。
\end{example}
非对易规范场称为 \textbf{Yang--Mills 场}。
式\eqref{chfb:eqn_GGamma}给出其局部规范势，式\eqref{chfb:eqn_GDG}给出规范变换。
把\eqref{chfb:eqn_GGamma}式代入曲率公式\eqref{chfb:eqn_R-vector-bundle}（$R_{\alpha i j}^\beta
= {\partial_i} \Gamma^{\beta}_{j \alpha} -\partial_j \Gamma^{\beta}_{i \alpha}
+\Gamma^{\sigma}_{j \alpha} \Gamma^{\beta}_{i \sigma}
-  \Gamma^{\sigma}_{i \alpha}\Gamma^{\beta}_{j \sigma}$），有：
\begin{align}
	R_{\alpha i j}^\beta
	=&{\partial_i} A^a_j (T_a)_{\alpha}^\beta
	-\partial_j A^a_i (T_a)_{\alpha}^\beta
	+A^c_j (T_c)_{\alpha}^\sigma  A^b_i (T_b)_{\sigma}^\beta
	-A^b_i (T_b)_{\alpha}^\sigma  A^c_j (T_c)_{\sigma}^\beta \notag\\
	=& \sum\nolimits_{a=1}^{r}(T_a)_{\alpha}^\beta  \times \bigl(
	{\partial_i} A^a_j  -\partial_j A^a_i  + \sum\nolimits_{b,c=1}^{r} A^b_i A^c_j C_{bc}^a    \bigr) .
	\label{chfb:eqn_YM-Strength}
\end{align}
其中$C_{bc}^a$是Lie 代数$\mathfrak{g}$的结构常数（注意$(T_a)_{\alpha}^\beta $是常数矩阵）。
式\eqref{chfb:eqn_YM-Strength}正是非对易规范场的场强公式，它是非对易规范场论的基础。
当$G$可对易时，\eqref{chfb:eqn_YM-Strength}式退化为\eqref{chfb:eqn_EM-strength}式。
因此，规范场强可视为相应联络的曲率。
\qed

\begin{remark}
    本节只给出规范势、规范变换与曲率的局部表达。
    全局构造还需指定 $G$-主丛、其联络及 $G$ 在纤维上的表示。
\end{remark}

\begin{remark}
    平直时空上的 Yang--Mills 理论也可直接用上述局部分量公式表述。
\end{remark}

